\programdayheader{Tuesday}
\daypalettebanner{Tuesday}

\section*{Morning}
\addcontentsline{toc}{section}{Tuesday -- Morning}

\subsection*{A-1502}
\textbf{Theme:} Gravitational waves / High energy astrophysics

\hypertarget{abstract-talk-283}{}\label{abstract-page-talk-283}
\subsubsection*{Using CFHT's SITELLE to probe the long-sought shell in the Crab nebula}
\addcontentsline{toc}{subsubsection}{Using CFHTs SITELLE to probe the long-sought shell in the Crab nebula}
\textbf{Speaker:} Lucas da Conceição\\
\textbf{Fields:} Compact Objects $\cdot$ Gas, dust and the ISM $\cdot$ High Energy Astrophysics

The Crab Nebula, the archetypal Pulsar Wind Nebula, presents a long-standing enigma: the absence of its expected supernova remnant shell. This missing component is critical for understanding the supernova's explosion mechanism and progenitor. We present a deep integral field spectroscopic study of the environment of the Crab Nebula using the Canada-France-Hawaii Telescope's SITELLE instrument, specifically targeting the [Fe XIV] $\lambda$5303 coronal line, a key tracer of hot, shocked gas, across two 11$'$ $\times$ 11$'$ fields. Our methodology consists of integral field unit (IFU) spectroscopy to construct detailed flux and velocity maps of the [Fe XIV] line outside of the visible nebula. Our primary objective is to derive plasma temperature estimates and evaluate whether the observed emission properties align with theoretical predictions for the forward shock. By combining observational data with advanced modelling techniques, we aim to shed light on the long-standing puzzle around the missing shell in this iconic pulsar wind nebula. In addition, our work highlights SITELLE's unique capability for wide-field, high-resolution searches of faint supernova remnant structures and provides crucial insights into the formation and evolution of shell-less Pulsar Wind Nebulae. This study contributes to a broader program aimed at understanding Pulsar Wind Nebulae across their evolutionary stages, including investigations of the dynamics and particle acceleration in systems such as CTB 80 (for which SITELLE data were awarded) and searches for continuous gravitational-wave signatures from PWN candidates.

\hypertarget{abstract-talk-13}{}\label{abstract-page-talk-13}
\subsubsection*{The First Uncontaminated Maps of Supernova Remnant Sgr A East}
\addcontentsline{toc}{subsubsection}{The First Uncontaminated Maps of Supernova Remnant Sgr A East}
\textbf{Speaker:} Mayura Balakrishnan\\
\textbf{Field:} High Energy Astrophysics

We present the first X-ray maps of supernova remnant (SNR) Sgr A East that are uncontaminated by emission from the surrounding Sgr A* environment, using Generalized Morphological Component Analysis (GMCA) applied to stacked Chandra ACIS-I observations of the Galactic Center. This approach enables, for the first time, a clean separation of the supernova remnant emission from the WolfRayet wind-fed plasma associated with Sgr A*, which overlap along the line of sight. By decomposing the X-ray emission and comparing with multiwavelength data, we identify the reflected shock and produce spatially resolved maps of Fe and S/Ar/Ca emission, finding centrally concentrated Fe consistent with mixed-morphology SNRs. Isolating the SNR emission in spectral modeling yields more robust plasma parameters: a lower ionization temperature and much higher electron density than previously reported, consistent with the interaction with the surrounding molecular material. Our results support dust survival through the reflected shock and indicate that Sgr A East exploded within a molecular filament accreting toward Sgr A*, vaporizing dense clumps and allowing the hot plasma to expand into low-density cavities, offering a rare, resolved view of supernova evolution in an extreme nuclear environment.

\hypertarget{abstract-talk-248}{}\label{abstract-page-talk-248}
\subsubsection*{Unifying the Neutron Star Zoo: Radio and X-ray Observations of Glitches in PSR J2229+6114}
\addcontentsline{toc}{subsubsection}{Unifying the Neutron Star Zoo: Radio and X-ray Observations of Glitches in PSR J2229+6114}
\textbf{Speaker:} Wenke Xia\\
\textbf{Field:} High Energy Astrophysics

Pulsars and magnetars are typically classified as two classes of neutron stars that differ in their inferred surface magnetic field (B-field) strengths, with distinct emission properties and activities. However, this dichotomous classification is challenged by the transient magnetar-like emission and bursting behaviour observed in two high-magnetic-field (high-B) pulsars immediately after large rotational glitches. Those two pulsars are believed to be transitional objects that bridge the emission and bursting behaviours of pulsars and magnetars. They lead to a neutron-star model in which the B-field strength serves as a unifying parameter. However, such a unifying picture remains incomplete due to the lack of constraints at the lower end of the B-field distribution. This is primarily because the rapid X-ray follow-up of pulsar glitches is rare, as the required near-daily glitch monitoring is generally not possible at most radio observatories. In this talk, we will present our first result from an ongoing near-real-time glitch monitoring campaign at the CHIME telescope. Using contemporaneous observations from the telescope's CHIME/Pulsar instrument (radio) and the NICER mission (X-ray) of low-B-field PSR J2229+6114, we identified three glitch events that occurred between 2020 and 2024 and analyzed the source's emission surrounding these glitches. Our results show no measurable glitch-induced changes in the pulsars X-ray emission, in stark contrast to the post-glitch activities in the emission of high-B pulsars. The results impose further constraints on the post-glitch behaviour of low-B-field pulsars, providing additional evidence that the glitch-induced emission changes observed in high-B pulsars are probably unique.

\hypertarget{abstract-talk-117}{}\label{abstract-page-talk-117}
\subsubsection*{Shining a Light on Ultraluminous X-ray Sources with HST and JWST}
\addcontentsline{toc}{subsubsection}{Shining a Light on Ultraluminous X-ray Sources with HST and JWST}
\textbf{Speaker:} Qiana Hunt\\
\textbf{Fields:} Compact Objects $\cdot$ High Energy Astrophysics $\cdot$ Transients

Ultraluminous X-ray sources (ULXs) are compact objectluminous star binaries with X-ray luminosities exceeding the standard Eddington limit for a 10 solar-mass black hole (Lx $>$ 10e39 erg/s). Given their extreme energetics, ULXs have been considered the prime candidates for finding intermediate-mass black holes, a rare and elusive category of black holes with masses between 100-100,000 solar masses. However, although nearly 2000 ULXs have been identified to date, no ULX has been found to definitively contain a black hole accretor, including the only the only known Galactic ULX. Thus, the nature of these high-energy sources remains a mystery. While characterizing the compact object component of ULXs is a complex issue, studying the stellar counterparts of ULXs is a crucial, accessible avenue of better understanding these binary systems. In this work, I conduct a multiwavelength survey of nearby extragalactic ULXs. Combining the high-resolution capabilities of the Chandra X-ray Observatory, the Hubble Space Telescope, and the James Webb Space Telescope, I directly identify the optical and near-infrared counterparts of 46 ULXs across 28 galaxies and perform a photometric and spectroscopic analysis of each, estimating the best-fit properties of all candidate donor stars. I present a first-look at these ULXs, many of which have never been observed before.

\hypertarget{abstract-talk-140}{}\label{abstract-page-talk-140}
\subsubsection*{Multi-messenger Astronomy with CanDIAPL: Identifying Kilonovae with Rubin and Gravitational Wave Alerts}
\addcontentsline{toc}{subsubsection}{Multi-messenger Astronomy with CanDIAPL: Identifying Kilonovae with Rubin and Gravitational Wave Alerts}
\textbf{Speaker:} Mohammed Chamma\\
\textbf{Fields:} Compact Objects $\cdot$ Gravitational waves $\cdot$ Transients

The upcoming Canadian Data-Intensive Astrophysics Platform (CanDIAPL) aims to build a powerful cloud computing centre for processing the enormous data streams that will come from facilities like the Vera C. Rubin Telescope and the Square Kilometer Array (SKA) precursors. An important component of CanDIAPL will be the Multi-Messenger Portal, a one-stop-shop for researchers to perform prompt follow-up and analysis of transient phenomena using data from multiple facilities and astrophysical messengers. Of particular interest is the identification of electromagnetic (EM) counterparts to gravitational wave (GW) events. We present here preliminary work for ingesting and processing the Rubin Observatorys alert stream from the Legacy Survey of Space and Time (LSST). This processing will classify alerts to identify potential EM counterparts to GW events detected by the LIGO-Virgo-KAGRA (LVK) observatories. The LSST alert stream is expected to generate 10 million alerts every night, or nearly 2 TB of alert data. LVK detections will have large localization regions on the sky that additionally make identifying a candidate counterpart difficult. These significant processing and scientific challenges require fine-tuned classification algorithms optimized for performance and accuracy. We discuss the deployment of custom software infrastructure on Canadian computing networks (such as CANFAR) to handle the processing and storage requirements. We also discuss light curve modelling, host association, and candidate filtering to identify kilonovae candidates. We aim to be processing alerts and searching for kilonova live in time for an upcoming limited LVK observing run slated to begin in Fall 2026.

\hypertarget{abstract-talk-134}{}\label{abstract-page-talk-134}
\subsubsection*{Are repeating and non-repeating FRBs one population? Insights from high-time resolution studies with CHIME/FRB}
\addcontentsline{toc}{subsubsection}{Are repeating and non-repeating FRBs one population? Insights from high-time resolution studies with CHIME/FRB}
\textbf{Speaker:} Alice Curtin\\
\textbf{Fields:} Compact Objects $\cdot$ High Energy Astrophysics $\cdot$ Transients

Fast Radio Bursts (FRBs) are a class of highly luminous, extragalactic radio transients that occur on nano-to-millisecond timescales. While most FRBs have been detected only once from a given sky location and distance, a small subset repeat, raising the question of whether there are multiple populations of FRBs. In this talk, I will investigate the repeating versus thus-far non-repeating FRB populations using the Canadian Hydrogen Intensity Mapping Experiment (CHIME)/FRB project. In particular, I will present the largest comparison of repeating and non-repeating FRBs at timescales down to microseconds, with over five years of data from CHIME/FRB. I will show that there are significant differences in the morphology (time and frequency structure) of the two populations, yet no large difference in their host galaxy environments. I will discuss the implications of this for a unified FRB population.

\subsection*{A-4502}
\textbf{Theme:} Stars and Stellar cluster formation

\hypertarget{abstract-talk-26}{}\label{abstract-page-talk-26}
\subsubsection*{LSPM J0207+3331: The first highly polluted cool hydrogen-rich white dwarf}
\addcontentsline{toc}{subsubsection}{LSPM J0207+3331: The first highly polluted cool hydrogen-rich white dwarf}
\textbf{Speaker:} Érika Le Bourdais\\
\textbf{Fields:} Exoplanets $\cdot$ Stars and Stellar cluster formation

Polluted white dwarfs serve as unique laboratories to study the bulk composition of exoplanetary material. As the remnants of the once complete planetary system orbit the dead star, parts can get trapped in the powerful gravitational field of the white dwarf and form a debris disk. The material will eventually disintegrate and accrete in the stars pristine hydrogen or helium atmosphere, polluting the white dwarf. While the field focuses mostly on warm (10,000 K $<$ Teff $<$ 25,000 K) helium-rich white dwarfs because the conditions are optimal for detecting metal pollution, I will report on the discovery of the first highly polluted cool hydrogen-rich white dwarf: LSPM J0207+3331. I will present various results ranging from the impact of metals on determining the physical parameters of cool white dwarfs in the near infrared to the bulk composition of the parent body polluting this white dwarf. This discovery demonstrates that planetary systems can maintain massive enough reservoirs of planetary material to sustain the pollution of a white dwarf billions of years after the death of its progenitor.

\hypertarget{abstract-talk-28}{}\label{abstract-page-talk-28}
\subsubsection*{Super-Earth Masses and Stellar Abundances from NIRPS Reveal Tentative Evidence for Water-Rich Formation around M Dwarfs}
\addcontentsline{toc}{subsubsection}{Super-Earth Masses and Stellar Abundances from NIRPS Reveal Tentative Evidence for Water-Rich Formation around M Dwarfs}
\textbf{Speaker:} Ryan Cloutier\\
\textbf{Fields:} Exoplanets $\cdot$ Stars and Stellar cluster formation

Tracing the compositional link between terrestrial super-Earths and their host stars provides clues of their dominant formation pathway. By constraining the stellar abundances of the refractory elements Mg, Si, and Fe, we can predict the core mass fractions (CMFs) of orbiting super-Earths. The level of agreement between this predicted CMF and the planets true CMF based on mass and radius measurements serves as a smoking gun of past formation processes such as mantle stripping and/or water-rich formation plus sequestration in the planets interior. I will present the first results from the Near Infrared Planet Searchers (NIRPS) CMF GTO program: an intensive radial velocity campaign to refine the masses of transiting super-Earths around M dwarfs and compute host stellar abundances. We retrieve precise masses for three hot super-Earths (GJ 1132 b, GJ 1252 b, and LTT 3780 b) of 1.69$\pm$0.15 (11$\sigma$), 1.54$\pm$0.18 (8$\sigma$), and 2.34$\pm$0.10 (23$\sigma$) Earth masses, respectively. We measure the CMFs of these and six additional hot super-Earths in the literature to 10-15\% precision, which when compared to CMF predictions from their host star abundances reveals that these planets are systematically lower density than is suggested by their host stars. This discrepancy suggests that hot super-Earths around M dwarfs habour significant reservoirs of water, generally consistent with water mass fractions of $\sim$1 - a few \%, which is likely sequestered in the interiors of these planets due to their high equilibrium temperatures.

\hypertarget{abstract-talk-76}{}\label{abstract-page-talk-76}
\subsubsection*{Investigating Jets and Winds from Three Time-Variable Protostars}
\addcontentsline{toc}{subsubsection}{Investigating Jets and Winds from Three Time-Variable Protostars}
\textbf{Speaker:} Doug Johnstone\\
\textbf{Field:} Stars and Stellar cluster formation

The formation of a star proceeds through the inside-out gravitational collapse of a protostellar core. Mass free-falls from the envelope toward the forming star, mediated by both magnetic fields and core-scale angular momentum, depositing material onto a rotationally supported disk. The Keplerian disk then feeds the central mass by reorganizing, or shedding, angular momentum. Finally, coupling between the magnetic field and disk rotation allows for powerful outflows to be driven from these systems, either as fast collimated jets launched close to the central protostar or as slower wide-angle winds from further out in the disk. In this talk I present the highlights from three recent investigations into the jets and winds from time-variable protostars, where the flow of mass through the disk and onto the star is episodic. From IRAS 04166+2706 (Tafalla et al. A\&A, 2025) we uncover with ALMA a remarkable series of arc-like features along the fast collimated jet, which are well fit as internal bow-shock working surfaces. Alternatively, for HOPS 358 (Kim et al. Nature Communications, 2025), ALMA reveals a much slower, wide-angle, rotating MHD-driven disk wind, efficiently extracting angular momentum and driving disk accretion. Finally, combining JWST with ALMA, we investigate the quasi-periodic protostars EC53 (Lee et al. Nature, 2025) during both burst and quiescence, witnessing the formation of crystalline silicates during episodic heating and spying the wide-angle disk wind responsible for launching these silicates to the outer reaches where they can become trapped in comets. Together, these three systems significantly advance our understanding of the coupling between star formation and stellar feedback.

\hypertarget{abstract-talk-70}{}\label{abstract-page-talk-70}
\subsubsection*{The Absolute Age of the Open Cluster NGC 6791 from Calibration Stars, Detached Eclipsing Binaries, and Synthetic CMD Fitting}
\addcontentsline{toc}{subsubsection}{The Absolute Age of the Open Cluster NGC 6791 from Calibration Stars, Detached Eclipsing Binaries, and Synthetic CMD Fitting}
\textbf{Speaker:} George Dufresne\\
\textbf{Field:} Stars and Stellar cluster formation

We determine an absolute age for the old, metal-rich open cluster NGC 6791 using Gaia DR3 photometry combined with detached eclipsing binaries (DEBs) in the cluster and three nearby Gaia calibration stars. We generate 10,000 Monte Carlo (MC) isochrone sets with DSEP, marginalizing over composition, convection and transport, microphysics, and key nuclear reaction rates, while sampling distance modulus and reddening over literature-motivated ranges. For each model we build a synthetic colormagnitude diagram (sCMD) that matches the observed star count in the MSTO and subgiant-branch window and injects empirical photometric scatter perpendicular to the ridgeline, enabling robust CMD comparisons without artificial-star tests. We evaluate CMD morphology with a bootstrap-calibrated 2D KolmogorovSmirnov statistic, and we incorporate external checks using nearest-point metrics: a coeval DEB statistic in (M,R,L) space and a non-coeval calibration-star statistic in (M\_G, BP-RP) space. These statistics are mapped to probability-density weights via Monte Carlo resampling and combined into a single isochrone weight. We infer an age of 8.42 $\pm$ 0.65 Gyr, [Fe/H]=+0.30 $\pm$ 0.10, Y=0.3006 $\pm$ 0.0201, (m-M) = 13.354 $\pm$ 0.05, and E(B-V)=0.177 $\pm$ 0.020. Our error budget shows no single dominant contributor, and highlights differences between open-cluster and globular-cluster age systematics. NGC 6791 remains a key benchmark for inner-disk chemical evolution, radial migration, and asteroseismic calibration.

\subsection*{A-3561}
\textbf{Theme:} Cosmology and Early Universe

\hypertarget{abstract-talk-293}{}\label{abstract-page-talk-293}
\subsubsection*{WebSky2.0: a pipeline for mock observations of the CMB, galaxies and LIM}
\addcontentsline{toc}{subsubsection}{WebSky2.0: a pipeline for mock observations of the CMB, galaxies and LIM}
\textbf{Speaker:} Nathan J. Carlson\\
\textbf{Field:} Cosmology and Early Universe

Upcoming cosmological surveys require mocks, simulated results that allow for direct comparison of observations to expected outcomes for various cosmological models. The WebSky simulations model the redshift-evolving distribution of dark matter halos, then apply response functions to reconstruct sky maps that are statistically analogous to observations. These simulations have been widely used in previous cosmic microwave background (CMB) research, but mass of resolutions are required for upcoming surveys. We present the WebSky2.0 mock CMB observations for the Simons Observatory (SO). Additionally, new WebSky2.0 tomographic line-intensity mapping (LIM) mocks shine a light into the formation of structure over cosmic time, and new WebSky2.0 mock galaxy catalogues provide new high resolution simulations of the low-redshift universe. This unified pipeline is ideal for comparative study between cosmological tracers, paving the way for future cross-correlation and oriented stacking works. WebSky maps can be generated from an arbitrary primordial \$\textbackslash{}zeta\$ field, opening up the search for traces of primordial physics in the large scale structure of the universe. Primordial intermittent non-gaussianities (PINGs) are of particular interest. PINGs are a generic feature that can arise during multi-field inflation, giving rise to isolated peaks in the matter field and may not be observable in conventional searches for primordial non-Gaussianity (PNG). We present the latest WebSky2.0 mocks for a suite of surveys and demonstrate their use in the search for PINGs and other signatures of new physics in the early universe. These include mocks of CMB observations with SO, galaxy surveys such as DESI and Euclid, and upcoming LIM surveys by the CCAT and COMAP collaborations.

\hypertarget{abstract-talk-69}{}\label{abstract-page-talk-69}
\subsubsection*{Searching for Primordial Wiggles: Testing Inflationary Particle Production with ACT DR6, Planck, and DESI DR2}
\addcontentsline{toc}{subsubsection}{Searching for Primordial Wiggles: Testing Inflationary Particle Production with ACT DR6, Planck, and DESI DR2}
\textbf{Speaker:} Simran Nerval\\
\textbf{Field:} Cosmology and Early Universe

The mechanism of inflation was developed in order to address problems that cannot be explained with the standard Hot Big Bang cosmological model, such as the spatial flatness and homogeneity of our universe. The simple picture of the inflationary epoch is one during which the universe underwent exponentially accelerating expansion within the first fractions of a second of the lifetime of the universe. The physical mechanism behind inflation is as yet unknown, but to date, we have been able to put an upper bound on the energy scale of this inflationary period using the Cosmic Microwave Background (CMB). The energy scale of inflation sets the temperature of the universe when inflation took place, which can be seen as a proxy for what physics was allowed during this time period. Our current upper bound on this energy scale requires models that go beyond the standard model of physics. One such exotic model includes inflationary particle production, which induces "wiggles", or oscillations, in the primordial density fluctuations that we see in the primordial power spectrum. I will present new constraints on these primordial wiggles and inflationary particle production using the Atacama Cosmology Telescope (ACT) Data Release 6 as well as Planck legacy data, baryon acoustic oscillation data from the Dark Energy Spectroscopic Instrument Data Release 2, and CMB lensing data from both ACT and Planck.

\hypertarget{abstract-talk-22}{}\label{abstract-page-talk-22}
\subsubsection*{The Glow of Axion Quark Nugget Dark Matter: CMB Spectral and Anisotropy Signatures}
\addcontentsline{toc}{subsubsection}{The Glow of Axion Quark Nugget Dark Matter: CMB Spectral and Anisotropy Signatures}
\textbf{Speaker:} Fereshteh Majidi\\
\textbf{Field:} Cosmology and Early Universe

Axion Quark Nuggets (AQNs) are macroscopic dark matter candidates formed during the QCD epoch through axion-induced charge separation. In this framework, dark matter consists of dense quark nuggets made of both matter and antimatter, naturally explaining the similarity between dark and visible matter abundances while behaving as cold dark matter on cosmological scales. A key feature of AQNs is that antimatter nuggets annihilate with ordinary baryons, producing a continuous injection of electromagnetic energy. In this talk, I present new predictions for the impact of AQN-induced energy injection on the cosmic microwave background (CMB). Using a modified version of the CLASS Boltzmann code, I compute the resulting $\mu$- and y-type CMB spectral distortions and examine their consistency with CMB anisotropy constraints. I find that AQNs produce a distinctive spectral distortion signal dominated by $\mu$-type distortions, with amplitudes $\mu \sim$ few $10^{-7}$ and $y \sim$ few $10^{-9}$. This hierarchy contrasts sharply with conventional decaying or annihilating particle dark matter models, where y-type distortions typically dominate. Despite the sizable spectral distortion signal, I show that AQN-induced heating has a negligible impact on CMB temperature and polarization anisotropies, with cosmological parameters remaining fully consistent with Planck $\Lambda$CDM constraints. These results demonstrate that CMB spectral distortions provide a powerful probe of macroscopic dark matter beyond the particle paradigm.

arXiv: 2512.05401

\hypertarget{abstract-talk-176}{}\label{abstract-page-talk-176}
\subsubsection*{Weak Gravitational Lensing Results from UNIONS}
\addcontentsline{toc}{subsubsection}{Weak Gravitational Lensing Results from UNIONS}
\textbf{Speaker:} Mike Hudson\\
\textbf{Field:} Cosmology and Early Universe

The Ultraviolet Near Infrared Optical Northern Survey (UNIONS) is a 6000 square degree imaging survey in the Northern Hemisphere, conducted at CFHT, Pan-STARRS and Subaru, for which weak lensing is a key science driver. The weak lensing team comprises members mostly from Canada and France, as well as Germany. I will review the weak lensing results from UNIONS, including its first measurements of the cosmological parameters, as well as new measurements of the density profiles of cosmic voids, intrinsic alignments, galaxy clusters, galaxy mergers and the correlation between galaxy halos and their black holes.

\section*{Afternoon}
\addcontentsline{toc}{section}{Tuesday -- Afternoon}

\subsection*{A-1502}
\textbf{Theme:} Next generation observatories - I

\hypertarget{abstract-talk-200}{}\label{abstract-page-talk-200}
\subsubsection*{Curing Cosmic Myopia: High-Resolution 10 MHz Imaging of Galactic Foregrounds from the High-Arctic with ALBATROS}
\addcontentsline{toc}{subsubsection}{Curing Cosmic Myopia: High-Resolution 10 MHz Imaging of Galactic Foregrounds from the High-Arctic with ALBATROS}
\textbf{Speaker:} Mohan Agrawal\\
\textbf{Fields:} Cosmology and Early Universe $\cdot$ Instrumentation $\cdot$ Next generation observatories

Accurate characterization of Galactic emission at decameter wavelengths is essential for detecting the faint 21-cm signal of neutral hydrogen from the cosmic dark ages. However, due to the challenges posed by human-made interference, and ionospheric effects, the state-of-the-art below 10 MHz still dates back to the 1960s, when Reber and Ellis caught brief glimpses of the Galaxy around 5 MHz with $\sim$5 resolution. The Array of Long Baseline Antennas for Taking Radio Observations from the Seventy-ninth Parallel (ALBATROS) is a new 8-antenna, 10-km-baseline drift-scan interferometer deployed on Axel Heiberg Island (79.8 N, 91.3 W) in the Canadian high Arctic that will image the northern sky between 520 MHz with arcminute resolution. An order-of-magnitude gain in resolution and coverage will pave the way for future dark ages experiments and open new avenues for studying ultra-low-frequency Galactic and extragalactic science. Because rugged terrain precludes real-time correlation, each ALBATROS station operates autonomously through the polar night and records baseband for offline correlation. Storage constraints require coarse channelization and 1-bit quantization, while power constraints necessitate low-power, high-jitter oscillators. New algorithms in ALBATROS GPU-accelerated data processing pipeline enable up-channelizing baseband to sub-Hz resolution, and tracking clock drifts to nanosecond accuracy using weather satellites. ALBATROS also has the capability to monitor the ionospheric plasma structure continuously, making solar and space-weather science possible. In this talk, I will provide an overview of the offline data-processing pipeline, and demonstrate early imaging and ionospheric science capabilities of the instrument using data collected during polar night.

\hypertarget{abstract-talk-335}{}\label{abstract-page-talk-335}
\subsubsection*{Nonlinear Spatiotemporal Wavefront Prediction for Adaptive Optics via Skip-Connected Koopman Autoencoders}
\addcontentsline{toc}{subsubsection}{Nonlinear Spatiotemporal Wavefront Prediction for Adaptive Optics via Skip-Connected Koopman Autoencoders}
\textbf{Speaker:} Artiom Matvei\\
\textbf{Fields:} Astrostatistics, Astroinformatics, and AI / Machine Learning $\cdot$ Exoplanets $\cdot$ Instrumentation

The next generation of ground-based astronomy will feature massive 30-to-40-meter mirrors designed to dramatically increase photon collection, enabling the observation of exceptionally dim targets such as exoplanets. However, because current large telescopes are fundamentally bottlenecked by atmospheric turbulence, they fail to reach their theoretical diffraction limits, meaning their increased scale is not yet fully realized. While Adaptive Optics (AO) systems correct for much of this turbulence, servo-lag and time-delay errors remain critical limitations. To fully unlock the potential of these giant mirrors and achieve the necessary performance boost for high-contrast imaging, advanced predictive control is essential to mitigate these delays and push resolutions closer to the diffraction limit. Traditional linear predictive filters, such as Empirical Orthogonal Functions (EOF) via truncated SVD, are widely used but can struggle with the highly non-linear, multi-scale nature of turbulent flow.

In this work, we present a novel deep learning approach for full-state wavefront prediction utilizing a Skip-connected Koopman Convolutional AutoEncoder (SKCAE). Adapted from complex fluid dynamics, this architecture uses a coordinate-transforming network to map high-dimensional, non-linear wavefront phases into a low-dimensional invariant subspace that evolves linearly. Furthermore, a Koopman dynamics-oriented subnetwork is utilized to quantitatively interpret the intrinsic mechanisms of the turbulence from a frequency perspective.

\hypertarget{abstract-talk-259}{}\label{abstract-page-talk-259}
\subsubsection*{Sharpening the Eyes of CASTOR}
\addcontentsline{toc}{subsubsection}{Sharpening the Eyes of CASTOR}
\textbf{Speaker:} Hanyu Xiao\\
\textbf{Field:} Instrumentation

CASTOR stands for the Cosmological Advanced Survey Telescope for Optical and ultraviolet Research. It is a Canada-led proposal for a next-generation space telescope. Ultraviolet light is of interest to astronomers for investigating the largest and brightest objects in the Universe, such as supermassive black holes. CASTOR's large field of view, equivalent to a full moon, allows for the study of celestial populations and rare transient events. CASTOR's capabilities make it a perfect complement to other modern space telescopes like the JWST and Euclid. CASTOR detector performance characterization is done at the Vacuum UltraViolet (VUV) Lab at the University of Calgary. The VUV Lab is one of the few Canadian vacuum UV laboratories capable of completing this work, as it can simulate the conditions in which the detectors will operate. I will go over features of the VUVLab at UCalgary that make it uniquely ideal for CASTOR's detector characterization, and highlights of the testing plan for the delta-doped CIS120 detectors, which have been treated by the Jet Propulsion Laboratory to increase UV sensitivity.

\subsection*{A-4502}
\textbf{Theme:} Transients

\hypertarget{abstract-talk-161}{}\label{abstract-page-talk-161}
\subsubsection*{Relativistic Maxwell-Bloch Equations with Applications to Astrophysics}
\addcontentsline{toc}{subsubsection}{Relativistic Maxwell-Bloch Equations with Applications to Astrophysics}
\textbf{Speaker:} Ningyan Fang\\
\textbf{Field:} Transients

We derive relativistic Maxwell-Bloch equations for potential applications in astronomical environments, where various radiative processes are known to occur, including the maser action and Dicke's superradiance. We show that for both phenomena a radiating system's response is preserved at different relative velocities between the system's rest frame and the observer, while the relevant timescales and the radiation intensity transform as expected from relativistic considerations. We verify that the level of coherence between groups of emitters travelling at different speeds is unchanged in all reference frames. We also derive relativistic versions of the maser equations applicable in the steady-state regime.

\hypertarget{abstract-talk-33}{}\label{abstract-page-talk-33}
\subsubsection*{Modelling the Energy and Redshift distribution of FRBs using the Second CHIME/FRB catalog}
\addcontentsline{toc}{subsubsection}{Modelling the Energy and Redshift distribution of FRBs using the Second CHIME/FRB catalog}
\textbf{Speaker:} Naman Jain\\
\textbf{Field:} Transients

Fast radio bursts (FRBs) are energetic millisecond-duration bursts with extragalactic origins. Roughly 4000 unique sources have been observed to date, about 3\% of which have been seen to repeat, though it remains unclear whether all FRBs repeat. An FRB population study addresses several key questions: How often do FRBs occur in the universe? What energy and redshift distributions do they follow? Is their occurrence correlated with the cosmic star formation history (SFH)? Answering these questions helps contextualize the origins of FRBs with potential astrophysical sources such as magnetars and supernovae. Previous population studies, based on samples of a few hundred sources, have provided important constraints on FRB rates and energetics. However, these analyses were limited by smaller sample sizes, less well-characterized selection effects, and uncertain treatment of repeat bursts. The recently released Second CHIME/FRB Catalog contains more than 90\% of all known FRB sources ($\sim$3600) and provides well-calibrated selection effects. We leverage this sample to infer key FRB properties, including the volumetric rate, energy function, redshift distribution, and correlation with the cosmic SFH. We then compare the predicted energetics and abundances with proposed FRB models and contrast the inferred redshift distributions with those of FRBs with known host galaxies.

\hypertarget{abstract-talk-304}{}\label{abstract-page-talk-304}
\subsubsection*{Under the Hood of an FRB Localization Machine: Lessons from Commissioning the CHIME Outriggers}
\addcontentsline{toc}{subsubsection}{Under the Hood of an FRB Localization Machine: Lessons from Commissioning the CHIME Outriggers}
\textbf{Speaker:} Gurman Sachdeva\\
\textbf{Fields:} Instrumentation $\cdot$ Transients

Fast radio bursts (FRBs) are bright, extragalactic, millisecond radio-frequency transients whose origins remain mysterious. Robustly discriminating between proposed progenitor models requires establishing a statistical sample (hundreds to thousands) of FRBs localized to tens of milliarcseconds precision. Over the next few years, such a sample will be obtained using very-long-baseline interferometry (VLBI) with the Canadian Hydrogen Intensity Mapping Experiments Outrigger telescopes. The recently constructed Outrigger telescope network is actively collecting data on-sky, with commissioning nearing completion. I will present an overview of the CHIME/FRB Outrigger project and results from the commissioning of the final Outrigger telescope. From design to on-sky performance, I will provide insight into the challenges of commissioning a new radio telescope and discuss key lessons learned. Finally, I will showcase the capabilities of the Outriggers by presenting an FRB localized to 13 pc precision in a galaxy 40 Mpc away.

\hypertarget{abstract-talk-269}{}\label{abstract-page-talk-269}
\subsubsection*{The Origin of Circumstellar Material and Its Imprint on Core Collapse Supernova Light Curves}
\addcontentsline{toc}{subsubsection}{The Origin of Circumstellar Material and Its Imprint on Core Collapse Supernova Light Curves}
\textbf{Speaker:} Nan Jiang\\
\textbf{Fields:} High Energy Astrophysics $\cdot$ Transients

Many hydrogen-rich core-collapse supernovae (Type II CCSNe) exhibit signatures of circumstellar material (CSM) interaction. However, the origin of this CSM, whether from steady, long-term stellar winds or short-term, intense mass-loss episodes shortly before core collapse remains poorly understood. In this talk, I will first present long-term, multi-band, high-cadence observations of KSP-SN-2023e from the KMTNet Supernova Program, spanning a baseline from four years pre-explosion to two years post-explosion. KSP-SN-2023e is a sub-luminous Type II CCSN featuring an unusually extended plateau phase (136 days), a rapid decline during the first 50 days post-peak, and a low $^{56}$Ni yield. Pre-explosion HST observations and the low $^{56}$Ni mass suggest the explosion originated from a low-mass red supergiant progenitor. Hydrodynamical modeling using SNEC indicates that \$\textbackslash{}sim 0.9 M\_\{\textbackslash{}odot\}\$ of CSM is required to produce the early excess emission. Crucially, our pre-explosion KSP observations strongly disfavor a scenario where this substantial CSM originated from outburst-like activity. Building on this, I will discuss the role of CSM in setting the phenomenological transition between Type II-P and II-L subclasses, using KSP-SN-2022c as a case study. Early light-curve fitting of KSP-SN-2022c with a single power-law places our first detection at roughly 15 minutes after shock breakout. This event represents a transitional case between the II-P and II-L subclasses, characterized by a rapid rise and decline, a luminous peak, a short plateau, and blueward motion in \$B-V\$ color. SNEC simulations reveal that the luminous peak and fast decline are driven by CSM interaction, while the short plateau is the result of envelope stripping.

\subsection*{A-3561}
\textbf{Theme:} Gas, dust and the ISM

\hypertarget{abstract-talk-289}{}\label{abstract-page-talk-289}
\subsubsection*{Developing PAH diagnostic tools to probe the interstellar medium in galaxies}
\addcontentsline{toc}{subsubsection}{Developing PAH diagnostic tools to probe the interstellar medium in galaxies}
\textbf{Speaker:} Erica Roscoe\\
\textbf{Field:} Gas, dust and the ISM

Polycyclic aromatic hydrocarbons (PAHs) are important tracers of the chemical and physical evolution of the interstellar medium (ISM), including high-energy environments such as star-forming regions in galaxies and protoplanetary disks. It is therefore vital to probe the relationship between ISM phases and the PAH populations within them. In this work, we use JWSTs NIRSpec and MIRI-MRS to investigate how PAH size varies across the photo-dissociation region (PDR) of the reflection nebula NGC 7023, using the well-established 11.2$\mu$m/3.3$\mu$m ratio as a size diagnostic. Because PAH size is regulated by the local radiation field, NGC 7023, which contains multiple H2 dissociation fronts, provides an ideal environment to examine how PAHs evolve from exposed atomic regions into shielded, denser gas. JWSTs unprecedented spatial and spectral resolution allows us to identify subtle shifts in the 11.2$\mu$m PAH band profile due to a carrier that behaves independently of PAHs and we attribute to very small grains (VSGs; Khan et al. 2025), previously undetectable with Spitzer. By removing this non-PAH contribution, for the first time we obtain a clean measurement of the 11.2$\mu$m PAH feature for use in the 11.2$\mu$m/3.3$\mu$m diagnostic. In this work, we find that smaller PAHs are progressively destroyed as the radiation field strengthens in the atomic region, leaving behind larger PAHs. We also find the smallest PAHs gather at the first H2 dissociation front, i.e. the transition from atomic to molecular gas, while the largest PAHs reside deeper in the molecular cloud at the farthest dissociation front. With JWST, we can accurately assess PAH photochemical evolution in NGC 7023 and develop an uncontaminated tool for probing high-intensity regions of the ISM.

\hypertarget{abstract-talk-113}{}\label{abstract-page-talk-113}
\subsubsection*{In-situ PAH formation in an O-rich Planetary Nebula: The Butterfly Nebula}
\addcontentsline{toc}{subsubsection}{In-situ PAH formation in an O-rich Planetary Nebula: The Butterfly Nebula}
\textbf{Speaker:} Nicholas Clark\\
\textbf{Field:} Gas, dust and the ISM

The James Webb Space Telescope (JWST) has revealed unexpected spatial variations in the mid-infrared emission of the Butterfly Nebula (NGC 6302), a Planetary Nebula that has been oxygen-rich throughout its entire evolution. The nebula surrounds a stellar remnant with the highest known effective temperature in our Galaxy, and this extreme radiation field drives a rich and complex chemistry in the surrounding gas and dust. Among the most surprising products of this chemistry is an abundance of Polycyclic Aromatic Hydrocarbons (PAHs) - large, planar carbon-based molecules whose formation in an oxygen-rich environment poses a fundamental challenge to current astrochemical models. PAHs are thought to form in the circumstellar envelopes of carbon stars, but this mechanism cannot account for their presence in the Butterfly Nebula, where the oxygen-rich composition places severe constraints on the available carbon chemistry. The strong and spatially varying PAH emission observed across the nebula instead points to in-situ formation and ongoing chemical processing of these molecules within the ionized nebular environment. Here we present an analysis of the spatial distribution of the PAH emission features, discuss the physical and chemical conditions that may enable their formation, and provide an estimate of the in-situ PAH formation rates. These results offer direct observational constraints on one particular formation channel for PAHs, contributing to a broader understanding of how these molecules arise across diverse astrophysical environments.

\hypertarget{abstract-talk-154}{}\label{abstract-page-talk-154}
\subsubsection*{Resolved Ionized Gas in M33: A Spectroscopic Catalog of Thousands of HII Regions from SIGNALS}
\addcontentsline{toc}{subsubsection}{Resolved Ionized Gas in M33: A Spectroscopic Catalog of Thousands of HII Regions from SIGNALS}
\textbf{Speaker:} Emma Jarvis\\
\textbf{Field:} Gas, dust and the ISM

Using SITELLE observations from the Star formation, Ionized Gas and Nebular Abundances Legacy Survey (SITELLE), we study thousands of HII regions and the surrounding diffuse ionized gas (DIG) in the nearby spiral galaxy M33. As one of the closest spiral galaxies, M33 is an ideal target to study the resolved structure and physical conditions of the ionized interstellar medium (ISM) across a galactic disk. Our mosaic of nine 11x11 fields provides the largest integral field spectroscopic view of the ionized gas in M33 to date, enabling a comparison of ionized regions across a range of galactic environments. The high spectral (R$\sim$5000) and spatial resolution ($\sim$2 pc) of the data allow us to resolve the internal structure of individual HII regions and their interaction with the surrounding ISM. From the SITELLE datacubes we measure emission-line fluxes (H$\alpha$, S[II], N[II], O[III], H$\beta$, and O[II]) and derive key physical properties including region sizes, luminosities, velocities, dust attenuation, and diagnostic line ratios. These line ratios are used to infer gas-phase metallicity, ionization parameter, and electron density of individual regions, enabling us to study environmental trends and metallicity gradients across the disk. We also investigate the relative contribution of compact HII regions and the DIG to the total ionized emission budget of the galaxy, providing constraints on ionizing photon propagation. The resulting catalog of spectroscopically characterized HII regions constitutes a valuable legacy dataset for studies of ionized nebulae and a reference for interpreting spectroscopic observations of star-forming galaxies in the nearby and distant universe.

\hypertarget{abstract-talk-260}{}\label{abstract-page-talk-260}
\subsubsection*{Numerical Simulations of the Streaming Instability in Protoplanetary Disks}
\addcontentsline{toc}{subsubsection}{Numerical Simulations of the Streaming Instability in Protoplanetary Disks}
\textbf{Speaker:} Benjamin Cheung\\
\textbf{Field:} Gas, dust and the ISM

In protoplanetary disks, dust and gas orbiting a newly formed star undergo gravitational collapse, with dust particles growing from micron-sized grains up to kilometer-sized planetesimals. However, collisional and aerodynamical effects contribute to a "meter barrier", inhibiting the efficient growth of dust grains at the centimeter and meter size scales. A promising mechanism to overcome this barrier is the streaming instability (SI), a fluid instability resulting from the dust experiencing aerodynamic drag as it moves through the surrounding gas. Streaming allows the local dust surface density to greatly increase, encouraging the formation of planetesimals. Furthermore, streaming is dependent on observable disk properties, such as dust grain size and the dust-to-gas surface density ratio. Thus, the necessary disk conditions for triggering efficient streaming, and whether streaming can serve as the starting point for planet formation in general, are of particular interest. To this end, numerical simulations are required to model the non-linear phase of the SI. We present results of 3D simulations of the SI in protoplanetary disks using the Athena hydrodynamics grid code. While other studies of the SI assume a single fixed grain size, we use a distribution of sizes consistent with grain growth predictions. This allows us to study more realistic disk environments and place better constraints on the conditions required to trigger efficient streaming.

\programdayheader{Wednesday}
\daypalettebanner{Wednesday}

\section*{Morning}
\addcontentsline{toc}{section}{Wednesday -- Morning}

\subsection*{A-1502}
\textbf{Theme:} Stars and Stellar cluster formation

\hypertarget{abstract-talk-326}{}\label{abstract-page-talk-326}
\subsubsection*{JWST reveals unexpected Carbon chemistry in NGC 6302}
\addcontentsline{toc}{subsubsection}{JWST reveals unexpected Carbon chemistry in NGC 6302}
\textbf{Speaker:} Charmi Bhatt\\
\textbf{Field:} Gas, dust and the ISM

Planetary nebulae (PNe, singular: PN) are chemical laboratories where circumstellar material undergoes significant physical and chemical processing before its injection into the interstellar medium. These objects harbour a wide variety of molecules, ranging from small, simple species such as OH, CO$_2$, and SO, to large complex molecules such as C$_{60}$ and C$_{70}$. Such environments are classified as either C-rich or O-rich based on the C/O ratio in the stellar photosphere, which governs subsequent chemistry during the PN phase. Yet $\sim$50\% of Bulge PNe exhibit dual chemistry (i.e. the simultaneous presence of O-rich and C-rich molecules), which non-equilibrium chemical models cannot fully reproduce. Recent JWST observations of the planetary nebula NGC 6302 shed new light on this problem. NGC 6302 is an O-rich, butterfly-shaped nebula with an exceptionally hot central star (T\_eff $\sim$220,000 K) surrounded by a thick dusty torus. These observations unexpectedly revealed these two carbon-bearing species: First, we detect CH$_3^+$, a key driver of organic chemistry in irradiated environments. Spatially resolved observations show that CH$_3^+$ is co-located with CO, H$_2$, HCO$^+$, and polycyclic aromatic hydrocarbons. Second, we detect CO$_2$ gas and ice within the dusty torus, where dense dust shields against intense UV radiation and enables ice chemistry. CO$_2$ ice can facilitate the formation of complex organic molecules. Strikingly, the CO$_2$ gas-to-ice ratio exceeds that of young stellar objects by more than an order of magnitude, indicating distinct ice formation or processing mechanisms in the PNe environment. Incorporating these species into astrochemical models represents a key step toward addressing the dual-chemistry challenge.

\hypertarget{abstract-talk-223}{}\label{abstract-page-talk-223}
\subsubsection*{Rotation Rates of Newly Identified Young Clusters in the Solar Neighbourhood}
\addcontentsline{toc}{subsubsection}{Rotation Rates of Newly Identified Young Clusters in the Solar Neighbourhood}
\textbf{Speaker:} Leslie Moranta\\
\textbf{Field:} Stars and Stellar cluster formation

Recent surveys have revealed several previously unrecognized young stellar clusters within $\sim$200 pc of the Sun, offering new opportunities to study stellar angular momentum evolution at ages between a few Myr and a few hundred Myr. Characterizing the rotation properties of these nearby populations is essential for understanding how stellar spin evolves during the early stages of stellar evolution.

We analyze TESS photometric time series for members of $\sim$100 young clusters in the solar neighbourhood, including several recently identified systems. Rotation periods are measured using a Gaussian Process framework designed to model quasi-periodic stellar variability while limiting overfitting through entropy-based regularization. This automated approach enables robust rotation measurements across large samples of stars with heterogeneous light-curve quality.

Applying this method to thousands of TESS light curves, we derive rotation periods for thousands of stars spanning a wide range of stellar masses. The resulting rotation distributions reveal clear mass-dependent sequences within clusters and highlight significant diversity in spin rates at young ages. Newly identified clusters show rotation patterns broadly consistent with established angular momentum evolution trends, while also extending empirical constraints in the poorly sampled $\sim$200--800\,Myr regime.

These measurements provide one of the largest homogeneous samples of stellar rotation periods for nearby young clusters. The results offer new benchmarks for models of stellar spin evolution and contribute to improving empirical rotationage relations in the solar neighbourhood.

\hypertarget{abstract-talk-121}{}\label{abstract-page-talk-121}
\subsubsection*{The rotational and magnetic properties of Polaris: Challenges to standard stellar models}
\addcontentsline{toc}{subsubsection}{The rotational and magnetic properties of Polaris: Challenges to standard stellar models}
\textbf{Speaker:} James Barron\\
\textbf{Field:} Stars and Stellar cluster formation

Polaris, the North Star, is an enigmatic object. It is the closest classical Cepheid to Earth, and exhibits unusual stellar and pulsation properties that are challenging to reconcile with stellar evolution models. Polaris is a hypothesized merger product and may be a rare example of a Cepheid crossing the Hertzsprung gap and instability strip for the first time, making it key to understanding the evolutionary paths of intermediate-mass stars. We report the rotational and magnetic properties of Polaris derived from a five-year spectropolarimetric magnetic monitoring campaign performed using ESPaDOnS at the Canada-France-Hawaii Telescope. We calculate circularly polarized Least-Squares Deconvolution profiles and measure the associated longitudinal magnetic field (Bz) at high precision. The Bz variation remained extraordinarily stable over the observing span with unsigned Bz strengths $<$1 G. From the Bz variability, we determine a stellar rotation period of 100.29(19) days. This is the first direct measurement of the rotation period for any classical Cepheid or yellow supergiant. We measure the equatorial rotation velocity to high precision and compare it to stellar evolution tracks to help constrain Polaris properties on the main sequence. We set a conservative upper limit on the stellar inclination angle and find a high probability of strong misalignment between the stellar rotation and orbital axes. The complexity of Polariss magnetic signatures contrasts with the long-term stability of the Bz measurements, making it difficult to confidently ascribe an origin to the magnetic field. We discuss these results in the context of our broader studies on the magnetic properties of classical Cepheids and other classes of evolved intermediate- and high-mass stars.

\hypertarget{abstract-talk-6}{}\label{abstract-page-talk-6}
\subsubsection*{Filamentary Hierarchies and Superbubbles: The role of galactic processes on filament formation}
\addcontentsline{toc}{subsubsection}{Filamentary Hierarchies and Superbubbles: The role of galactic processes on filament formation}
\textbf{Speaker:} Rachel Pillsworth\\
\textbf{Field:} Gas, dust and the ISM

Large scale phenomena in spiral galaxies such as shear, supernovae and magnetic fields contribute to the formation and evolution of filamentary structures in their disks. While recent high resolution observations using JWST and ALMA have begun characterizing properties of both superbubbles and filaments across the disks of nearby galaxies, simulations offer important insight on the largest, kpc-scale filamentary structures the first steps in a multiscale star formation process. These structures are important to understanding the ISM in a galactic disk. I will present work from two recent papers which highlight the formation and evolution of kpc filaments in a galactic disk. Using filament and superbubble structure analysis codes, we find that magnetic fields play a significant role in the stability of a filament while galactic shear plays a smaller one. Additionally, we introduce a new parameter for filament stability, utilizing the surface and central pressures of filaments to distinguish between formation and fragmentation. Finally, I will compare these results over the evolution of the filament populations in both a Milky Way-like and an NGC 628-like isolated galactic disk simulation. By analyzing the filamentary hierarchy in each disk over time, we highlight the roles of accretion and spiral arms in maintaining filaments as long-lived structures. By comparing the filaments within spiral arms with those in interarm or flocculent regions, we further pinpoint the role of the overall galactic environment on the star formation in its disk.

\hypertarget{abstract-talk-328}{}\label{abstract-page-talk-328}
\subsubsection*{ALMA-SAILS: An Archival Search for Protostellar Accretion Streamers - The initial look into Serpens Main}
\addcontentsline{toc}{subsubsection}{ALMA-SAILS: An Archival Search for Protostellar Accretion Streamers - The initial look into Serpens Main}
\textbf{Speaker:} Samuel Fielder\\
\textbf{Field:} Gas, dust and the ISM

Protostellar accretion streamers are coherent filamentary structures that channel material from envelope scales and beyond into protostellar cores. This phenomena has emerged as an important yet poorly characterized mode of mass assembly during the early stages of star formation. While accretion streamers have long been predicted by simulations of protostellar collapse, they have only recently begun to appear with prevalence in observations, leaving their occurrence rates and physical properties across different star-forming environments largely unconstrained. Here we present initial results from ALMA-SAILS (ALMA Search for Asymmetry and Infall at Large Scales), a systematic survey exploiting the ALMA archive to characterize the incidence and morphology of accretion streamers in nearby star-forming regions. We focus on the first target region, Serpens Main, and report continuum-based statistics on the detected streamer candidate population, including occurrence rates as a function of protostellar properties. These results provide a first census of streamer incidence in this region and lay the groundwork for follow-up line observation analysis aimed at confirming infall kinematic signatures.

\hypertarget{abstract-talk-164}{}\label{abstract-page-talk-164}
\subsubsection*{FRBs, Pulsars, and AU-Scale Plasma: Probing the Universe with Scattering Screens}
\addcontentsline{toc}{subsubsection}{FRBs, Pulsars, and AU-Scale Plasma: Probing the Universe with Scattering Screens}
\textbf{Speaker:} Sachin Pradeep Etakkepravan Thulicheri\\
\textbf{Fields:} Compact Objects $\cdot$ Gas, dust and the ISM $\cdot$ Transients

Fast Radio Bursts (FRBs) and pulsars emit bright, millisecond-duration radio pulses that are scattered by AU-scale plasma structures along the line of sight. Scattering leads to multi-path propagation and temporal broadening, while scintillation the radio analogue of optical twinkling arises from coherent interference between these paths, imprinting frequency-dependent spectral modulation. Scintillation can place strong constraints on the size of emission regions in a variety of transients including FRBs, pulsars and GRB afterglows. However, robust interpretation of observations requires a theoretical and simulation framework, which we present in S Pradeep E. T. et al. 2025. In this framework, the suppression or absence of scintillation indicates decoherence caused by a scattering screen that spatially resolves the emission region. We make quantitative predictions for scintillation observables to detect such systems and constrain the size of the transient emission region and its frequency evolution. Pulse from a point source encountering two screens that do not spatially resolve each other shows two scintillation scales. When one screen is instead resolved by the other, one of those scales is suppressed, producing a distinctive signature in the frequency evolution of scintillation. For FRBs, which often encounter both a Galactic and a host-galaxy scattering screen, our method turns scintillation into a precision probe of AU-scale structures and turbulence in extragalactic plasma. When an intervening galaxy lies along the sightline, these measurements can directly reveal the clumpy structure of its circumgalactic medium. Through this work, we present a new tool to probe transient emission physics and the line-of-sight plasma structures.

\hypertarget{abstract-talk-101}{}\label{abstract-page-talk-101}
\subsubsection*{A SITELLE catalog of supernova remnants in nearby dwarf galaxies of the SIGNALS survey}
\addcontentsline{toc}{subsubsection}{A SITELLE catalog of supernova remnants in nearby dwarf galaxies of the SIGNALS survey}
\textbf{Speaker:} Pénélope Glasman\\
\textbf{Fields:} Galaxy $\cdot$ Gas, dust and the ISM

We present a comprehensive study of supernova remnants (SNR) in nine nearby dwarf galaxies observed with the imaging Fourier transform spectrometer SITELLE at the Canada-France-Hawaii Telescope. SITELLEs ability to map nebular emission with high velocity precision and to disentangle complex kinematic components provides a powerful new approach for identifying and characterizing SNR in external galaxies. This homogeneous dataset represents the largest and most uniform SITELLE-based SNR survey for dwarf galaxies. Using spatially resolved emission-line ratios and detailed kinematic diagnostics, we reassess all previously reported SNR candidates and identify numerous new remnants across several galaxies. The resulting catalog provides a detailed view of SNR kinematics, such as expansion velocities, shock properties, and local ISM conditions. Together, these measurements offer a coherent set of indicators that help place the remnants along their evolutionary sequence. In addition to the identification of individual remnants, this work offers new insight into the origin and morphology of diffuse ionized gas in dwarf galaxies, a key open question in ISM physics. Our results also provide a quantitative framework for evaluating the detectability of SNR in the optical regime and for assessing their role in the chemical enrichment and energy injection in low mass galaxies.

\hypertarget{abstract-talk-124}{}\label{abstract-page-talk-124}
\subsubsection*{Disentangling Observations of the Galactic Center with Score-Based Diffusion Models}
\addcontentsline{toc}{subsubsection}{Disentangling Observations of the Galactic Center with Score-Based Diffusion Models}
\textbf{Speaker:} Zach Sumners\\
\textbf{Fields:} Astrostatistics, Astroinformatics, and AI / Machine Learning $\cdot$ Gas, dust and the ISM $\cdot$ Supermassive black hole

The Galactic Center is a dense and dynamic region that makes studying the black hole at the center of the Milky Way, Sagittarius A* (Sgr A*), challenging especially in highly extincted regimes such as the mid-infrared. In this work, we develop a machine learning framework using score-based diffusion models to disentangle extended galactic structures, instrumental noise, and point-like signal for independent study. This data-driven technique learns unique morphological information of each structure type, then uses a Bayesian framework to generate separated images from a single input. It operates on the data-level with no parametric assumptions, flexibly encoding all information through training data. As a result, the technique is broadly applicable to a range of problems that require separating morphologically unique, but complex, components. We test the framework with mid-infrared images of Sgr A* from the James Webb Space Telescope, and recover visually and statistically plausible separated components consistent with known Galactic Center information. As the nearest supermassive black hole, Sgr A* remains an important observational target for understanding low-luminosity accretion, and continued development of signal separation techniques will further enhance the reliability of precision measurements.

\subsection*{A-4502}
\textbf{Theme:} Galaxy

\hypertarget{abstract-talk-316}{}\label{abstract-page-talk-316}
\subsubsection*{Thermodynamic Properties of Cool Core and NonCool Core Galaxy Cluster Atmospheres}
\addcontentsline{toc}{subsubsection}{Thermodynamic Properties of Cool Core and NonCool Core Galaxy Cluster Atmospheres}
\textbf{Speaker:} Marie-Joëlle Gingras\\
\textbf{Fields:} Galaxy $\cdot$ Galaxy Clusters

Galaxy clusters are the largest gravitationally-bound and virialized structures in the Universe. Most of their baryons are not in stars, but in a hot, X-ray emitting atmosphere known as the intracluster medium (ICM). This plasma retains the history of the long-term interplay between cluster assembly, radiative cooling, and the energetic feedback from supermassive black holes. X-ray observations have shown that the thermodynamics of the ICM are tightly linked to cooling, star formation, and AGN activity in galaxy clusters, giving rise to the long-standing cool core (CC) / non-cool core (NCC) dichotomy. While this dichotomy is well-known, the underlying physics for this separation are not yet fully understood.

In this talk, I will present a comprehensive analysis of the thermodynamic profiles of the hot X-ray atmospheres of 85 galaxy groups and clusters drawn from the Chandra archive. I will discuss commonly used CC diagnostics, testing them against multi-wavelength indicators of cooling and AGN activity. I will show how far CC and NCC differences extend for various thermodynamic properties, and how these differences become clearer once mass-driven scatter is reduced by scaling all clusters to a common mass, persisting to \$\textbackslash{}sim 0.7-1\$ R\$\_\{2500\}\$ for most properties. I will also present the median CC and NCC thermodynamic profiles obtained by combining the clusters in our sample, and compare them to hydrodynamical simulations. These median profiles also show that cool cores have accumulated a substantial excess of hot gas in their centres, corresponding to 1.5-2.5\% of \$M\_\{2500\}\$ compared to NCCs, and I will explore possible origins of this excess and what it may imply for the growth of cooling cores.

\hypertarget{abstract-talk-160}{}\label{abstract-page-talk-160}
\subsubsection*{Parsec Scale Analysis of the PAHs and Star Formation of Evil Eye Galaxy (M64)}
\addcontentsline{toc}{subsubsection}{Parsec Scale Analysis of the PAHs and Star Formation of Evil Eye Galaxy (M64)}
\textbf{Speaker:} Woorak Choi\\
\textbf{Fields:} Galaxy $\cdot$ Gas, dust and the ISM

The physical state of the interstellar medium (ISM) is deeply intertwined with a galaxy's evolutionary history and localized star formation activity. We present JWST NIRCam and MIRI observations of M64, the "Evil Eye" galaxy, to investigate the multi-phase ISM in a known post-minor merger system. Combining these data with ALMA CO(3-2) observations, we conduct a parsec-scale analysis of the interplay between dust, molecular gas, and stellar feedback. We derive continuum-subtracted maps of relatively dust-free star formation tracers (Pa\$\textbackslash{}alpha\$, Br\$\textbackslash{}alpha\$) and multiple PAHs (3.3, 7.7, and 11.3 \$\textbackslash{}mu\$m) to analyze the spatial distribution of star-forming regions and ISM/dust properties. We will present spatially resolved scaling relations between PAHs, ionized gas, and molecular gas across the inner disk (r $<$ 1 kpc) of M64. Our discussion will focus on utilizing PAH ratios and the overall PAH fraction as diagnostics for the local radiation field, gas density, PAH size distribution and their characteristics and the impact of stellar feedback on the ISM. Furthermore, we will explore how the minor merger affects the star formation and overall properties of the ISM.

\hypertarget{abstract-talk-187}{}\label{abstract-page-talk-187}
\subsubsection*{A new, realistic model for radiation and molecular gas in simulations of galaxies}
\addcontentsline{toc}{subsubsection}{A new, realistic model for radiation and molecular gas in simulations of galaxies}
\textbf{Speaker:} Padraic Odesse\\
\textbf{Fields:} Galaxy $\cdot$ Gas, dust and the ISM

Radiation is the dominant factor for establishing the multiphase structure of the interstellar medium (ISM). Extreme-ultraviolet radiation from young star clusters carves out hot HII regions, while far-ultraviolet radiation permeates further into photodissociation regions (PDRs) to destroy molecular gases and deliver heating to the cool ISM. All of these processes can occur at very small scales, with HII regions and PDRs having sizes around $\sim$1-10 pc, but even the best galaxy simulations cannot hope to resolve structures on this scale! Individual resolution elements in galaxy-scale simulations can be extremely optically thick ($\tau$ $\sim$10 to 1e5), so the radiation field can change drastically over just one particle or cell. All simulations must therefore include some correction to distribute radiation over a whole resolution element. Such corrections have been in use for over 20 years (e.g. Mellema et al. 2006), but we have improved on these works by adding a molecular component, and have laid the groundwork for a unified model that captures both HII regions and PDRs. In this talk, I will discuss the novel framework that I have developed for radiation on the subgrid scale. I will present new galaxy simulations that use this radiative transfer framework to produce realistic distributions of molecular gas, and make direct comparisons to observed atomic and molecular surface density profiles of galaxies from the THINGS survey. This model is necessary to fully capture the interplay between radiation and molecular gas in the dense ISM, which brings our simulations up-to-date with recent observational advances from surveys like PHANGS that deliver high-resolution gas data.

\hypertarget{abstract-talk-222}{}\label{abstract-page-talk-222}
\subsubsection*{KILOGAS: A Centrally Peaked but Spatially Extended Enhancement of Molecular Gas in Galaxy Post-Mergers}
\addcontentsline{toc}{subsubsection}{KILOGAS: A Centrally Peaked but Spatially Extended Enhancement of Molecular Gas in Galaxy Post-Mergers}
\textbf{Speaker:} Blake Ledger\\
\textbf{Field:} Galaxy

Simulations predict that the gravitational torques induced by the interaction of two merging galaxies can trigger radial inflows and funnel gas inwards. To date, observations have only been able to confirm this idea indirectly: galaxy mergers show centrally enhanced star formation rates (SFRs), low gas-phase metallicities, and increased supermassive black hole accretion rates, which imply an increase in available gas. Direct confirmation of the distribution of gas requires resolved observations in a large and representative sample of galaxy mergers with a homogeneously-observed sample of non-merger controls, but previous studies have been limited to a handful of individual systems or small samples.

Our goal is to directly measure the radial distribution of the molecular gas (i.e. the fuel for star formation) in galaxies that are in the final stage of the merger sequence: post-mergers. We use new kpc-scale CO (2-1) observations from the KILOGAS survey, which targets a representative sample of $\sim$500 galaxies with resolution-matched ALMA and optical IFU spectroscopy. From the optical imaging, we visually identify $\sim$50 post-mergers and a control pool of $\sim$200 non-interacting galaxies in KILOGAS. We find that the **global** molecular gas content of this population of post-mergers is comparable to that of the population of matched non-merger controls. However, on **resolved** scales, the post-mergers show a centrally peaked and spatially extended (out to $\sim$3 kpc) enhancement of molecular gas compared to the controls. Our results provide direct observational evidence for radial inflows of molecular gas in a statistically significant sample of post-mergers.

\hypertarget{abstract-talk-136}{}\label{abstract-page-talk-136}
\subsubsection*{The Photometric \& Overdensity Analysis of Titans in the Early Universe}
\addcontentsline{toc}{subsubsection}{The Photometric \& Overdensity Analysis of Titans in the Early Universe}
\textbf{Speaker:} Joe Bhangal\\
\textbf{Fields:} Galaxy $\cdot$ Galaxy Clusters

Large-scale structure in the early Universe remains poorly constrained, as simulations struggle to reproduce the most extreme overdensities observed at high redshift. Protoclusters, the progenitors of present-day galaxy clusters, are key testbeds for studying how galaxies and their environments co-evolve during the early stages of structure formation. Massive galaxies have long been proposed as signposts of protocluster environments due to their association with massive dark matter halos. We investigate whether 11 massive galaxies ($\geq 10^{12}$ Msun dynamical mass) trace overdense environments at redshift z$\sim$4-5. Deep JWST NIRCam/MIRI and Hubble imaging are used to identify overdensities via photometric redshift estimates. Of the 11 massive galaxies, we find 8 reside in significant overdensities at z$\sim$4-5, each containing on average $\sim$40 galaxies. The stellar mass function constructed from these environments shows a significant excess of massive galaxies relative to the field, with a flattening of the high-mass end instead of the typical exponential decline. The galaxies in the overdensities have star formation rates consistent with the field, suggesting the excess stellar mass results from an increased number of galaxies rather than enhanced star formation. These results suggest that massive galaxies are effective signposts of overdense environments at z$\sim$4-5 and that their environments may capture the early growing phase of cluster formation. Identifying such systems allows us to investigate the physical processes driving early cluster assembly, including the roles of star formation, feedback, and environmental effects in shaping galaxy evolution.

\hypertarget{abstract-talk-178}{}\label{abstract-page-talk-178}
\subsubsection*{Sweating the Small Stuff: Dwarf galaxy quenched fractions out to cosmic noon}
\addcontentsline{toc}{subsubsection}{Sweating the Small Stuff: Dwarf galaxy quenched fractions out to cosmic noon}
\textbf{Speaker:} Ian Roberts\\
\textbf{Field:} Galaxy

The quenching of star formation in galaxies is a fundamental component of galaxy evolution. At z$\sim$0, the current paradigm is one where the fraction of quenched galaxies evolves strongly with galaxy mass. The vast majority of high-mass galaxies (log M$_\star$ $\sim$11) are quenched and the vast majority of low-mass galaxies (log Mstar $\sim$9) are star-forming, however due to depth limitations and incompleteness, there are comparatively few constraints on galaxy quenching in the dwarf galaxy regime (log Mstar $<$ 9). The constraints that do exist come from the local group and nearby Milky Way analogs. These works find evidence for an increase in quenched fraction in the dwarf galaxy regime, suggesting that the efficiency of galaxy quenching is minimized around a critical stellar mass of $10^{9}$ Msun. Despite these intriguing results, there are no constraints on quenched fractions for these low-mass galaxies outside of the local Universe. Here we present the first such measurements. Using deep JWST photometry from COSMOS-Web, we measure galaxy quenched fractions down masses of $10^{7}$ - $10^{8}$ Msun, from the low-redshift Universe all the way out to Cosmic Noon (z $\sim$2.5). We measure an upturn in galaxy quenched fractions for stellar masses below $10^{9}$ Msun, the strength of which is strongest at z $\sim$0 and becomes shallower moving back in cosmic time. We postulate that this time evolution is related to the increasing importance of environmental quenching at low redshift, and present preliminary constraints on this topic. Low-mass galaxies are the most numerous in the Universe. Fully comprehending galaxy evolution requires understanding how these galaxies evolve and form stars across cosmic time.

\subsection*{A-3561}
\textbf{Theme:} Exoplanets - I

\hypertarget{abstract-talk-233}{}\label{abstract-page-talk-233}
\subsubsection*{Constraining the multiplicity of the coldest brown dwarfs down to planetary masses between 0.5 and 1000 au with JWST}
\addcontentsline{toc}{subsubsection}{Constraining the multiplicity of the coldest brown dwarfs down to planetary masses between 0.5 and 1000 au with JWST}
\textbf{Speaker:} Thomas Vandal\\
\textbf{Fields:} Exoplanets $\cdot$ Stars and Stellar cluster formation

The James Webb Space Telescope (JWST) provides an unprecedented opportunity to study the coldest, lowest-mass and faintest products of star formation: Y-type brown dwarfs, whose masses can be as low as those of directly imaged giant exoplanets. Here, we leverage JWSTs NIRISS and NIRCam instruments to probe the multiplicity of 22 Y dwarfs. We analyze our observations with Kernel Phase Imaging (KPI), a generalization of interferometric techniques to regular images, which enables us to search for companions below the diffraction limit of the telescope. Thanks to JWSTs exquisite sensitivity at 4.8 um, our KPI analysis has the ability to detect companions with sub-Jupiter masses (mass ratios down to 0.1), and separations as low as 0.5 au. We use KPI to recover the known Y dwarf binary WISE 0336 AB, whose two components may have planetary mases. Since KPI has a limited field of view ($\sim$2 arcsec), we also look for companions at wider separations with traditional PSF injection recovery. Despite not providing new detections, this wide-separation search significantly extends the parameter space we are able to probe, unlocking separations up to 1000 au. As a whole, our analysis is sensitive to planetary-mass companions between 0.5 and 1000 au. This provides the most robust constraint to date on the multiplicity of Y dwarfs and represents an important step towards a complete understanding of star formation processes down to the lowest masses, which is crucial to properly understand the interplay between stellar and planetary formation.

\hypertarget{abstract-talk-224}{}\label{abstract-page-talk-224}
\subsubsection*{A Sub-Neptune Atmospheric Atlas: Toward Population-Level Trends with JWST}
\addcontentsline{toc}{subsubsection}{A Sub-Neptune Atmospheric Atlas: Toward Population-Level Trends with JWST}
\textbf{Speaker:} Amélie Gressier\\
\textbf{Field:} Exoplanets

Sub-Neptunes and super-Earths, with sizes between Earth and Neptune, are among the most common exoplanets in the Galaxy, yet they have no counterparts in our Solar System. If such a planet had formed between the rocky inner planets and the ice giants of our Solar-System, what would its atmosphere look like? Studying the atmospheres of these worlds provides a unique opportunity to probe planetary formation and volatile accretion, and to inform our understanding of planet formation. Recent JWST observations are beginning to reveal that many sub-Neptunes possess atmospheres enriched in heavy elements compared to the hydrogen-dominated atmospheres of giant planets such as Jupiter and Saturn. This enrichment implies a substantial fraction of volatiles mixed within their envelopes and suggests that atmospheric mean molecular weight may vary across the sub-Neptune population. One possibility is that lower-mass planets retain proportionally larger volatile fractions, potentially producing a trend between atmospheric mean molecular weight and planetary mass. In this talk, I will review current atmospheric constraints for sub-Neptunes from JWST transmission spectroscopy and place them in the broader context of planetary mass, temperature, and bulk composition. I will discuss how assembling a catalogue of sub-Neptune spectra could enable comparative population studies, allowing us to investigate trends in atmospheric enrichment and better understand the formation and evolution of these planets.

\hypertarget{abstract-talk-82}{}\label{abstract-page-talk-82}
\subsubsection*{Atmospheric compositions of hot Jupiters as tracers of planet formation history}
\addcontentsline{toc}{subsubsection}{Atmospheric compositions of hot Jupiters as tracers of planet formation history}
\textbf{Speaker:} Michael Balogh\\
\textbf{Field:} Exoplanets

Atmospheric abundances of giant exoplanets are increasingly used to infer their formation pathways, motivating upcoming population studies with facilities such as the ESA Ariel mission. However, the extent to which such inferences depend on assumptions about protoplanetary disk physics remains poorly understood. We present a population synthesis study of giant planet formation that combines pebble and planetesimal accretion with a simple equilibrium C/O chemistry model, and contrast disks dominated by viscous or magnetically-driven wind accretion. The inclusion of planetesimals in the wind--driven disks leads to more massive planets and a population of hot Jupiters that closely match the observed mass-metallicity sequence. The atmospheric ratio of C/O does not generally correlate well with initial planet location. Instead, it primarily reflects the timing and location of envelope enrichment relative to migration and the onset of runaway gas accretion. Bulk abundance ratios anchored to refractory elements (e.g. O/Si, C/Si) provide useful diagnostics of solid enrichment history, and reveal population-level differences between giant planets that end up within or beyond $>$1AU after 3\,Myr.

\hypertarget{abstract-talk-234}{}\label{abstract-page-talk-234}
\subsubsection*{Three years of atmospheric characterization with The Near Infra-Red Planet Searcher (NIRPS)}
\addcontentsline{toc}{subsubsection}{Three years of atmospheric characterization with The Near Infra-Red Planet Searcher (NIRPS)}
\textbf{Speaker:} Romain Allart\\
\textbf{Field:} Exoplanets

NIRPS is a fiber-fed high-resolution spectrograph assisted by adaptive optics installed on the 3.6m telescope of ESO at LaSilla, Chile. Operated simultaneously with HARPS, NIRPS covers the Y, J, and H bands. In this talk, I will describe the objectives of its consortium that has been allocated 725 GTO nights over 5 years. A third of this time is dedicated to in-depth spectral characterizations, providing detailed, high-fidelity, high-signal-to-noise transmission and emission spectra in preparation for the ELT, as well as large, comprehensive atmospheric and orbital architecture surveys, with the goal of observing over 70 exoplanets, ranging from ultra-hot Jupiters to temperate terrestrial planets. Over the program's first half, we have prioritized the best exoplanets across the exoplanet population, including many JWST targets. During the talk, I will highlight key results we have obtained since the start of operations, including detections of escaping atmospheres, atmospheric variability, and constraints on the atmospheric composition of warm Neptunes and hot Jupiters. I will then highlight a first statistical analysis of the atmospheric escape for a sample of 40 exoplanets we observed. These results will showcase the potential of NIRPS to deliver high-fidelity atmospheric spectra to constrain the formation and evolution of exoplanets at the statistical level. Looking ahead, I will discuss how extending NIRPS into the K band will unlock a new realm of atmospheric science, enhancing and securing long-term follow-up of exoplanet atmospheres to study their variability. Finally, NIRPS lays the groundwork for ELT spectrographs such as ANDES, and I will showcase how the lessons learned from NIRPS can be applied to the next generation of state-of-the-art instruments.

\hypertarget{abstract-talk-228}{}\label{abstract-page-talk-228}
\subsubsection*{Pandora the Explorer: Early Results and Mission Status}
\addcontentsline{toc}{subsubsection}{Pandora the Explorer: Early Results and Mission Status}
\textbf{Speaker:} Jason Rowe\\
\textbf{Field:} Exoplanets

The Pandora SmallSat Mission, a NASA Astrophysics Pioneers project, aims to provide a critical understanding of exoplanetary atmospheres by observing transits across a broad wavelength range, through transmission spectroscopy.A primary challenge in transmission spectroscopy is Stellar Contamination, where stellar heterogeneities (spots and faculae) can mimic or obscure planetary atmospheric signatures. Pandoras simultaneous visible and IR channels aim to correlate the spectral features with photometric variability, which will determine the stellar contamination in the transit signal. In this presentation, I will provide an update on the Pandora mission status and performance. Furthermore, I discuss the availability and quality of observational data collected to date and how Pandora may be beneficial to your further science goals. This includes an overview of science targets and various auxiliary science targets. These targets are diverse and include an exciting range of astrophysical objects, such as various types of exotic stars, including a large sample of chemically peculiar stars

\hypertarget{abstract-talk-284}{}\label{abstract-page-talk-284}
\subsubsection*{New Exoplanets and Demographics Trends Detected with Gaia DR3}
\addcontentsline{toc}{subsubsection}{New Exoplanets and Demographics Trends Detected with Gaia DR3}
\textbf{Speaker:} William Thompson\\
\textbf{Field:} Exoplanets

The astrometry technique for detecting exoplanets has been attempted since the 1940s, but only recently have the very first independent detections of planets been made. This technique, which detects planets via the wobble they induce on stars' positions in the night sky, is poised to deliver up to thousands of new exoplanets once Gaia DR4 is released late this year. Ahead of this major release, we present a novel technique for extracting independent exoplanet detections from *existing* Gaia DR2, DR3, and Hipparcos data. We will present of order $\sim$100s of new planets detected through this technique around nearby stars, thousands of candidate planets, and preliminary demographics trends inferred from this new large, uniform sample. In addition, our public catalogue will contain full Bayesian posteriors for companions, and non-detections (upper limits) for over 100,000 nearby stars. This catalogue can form the basis of independent demographics studies for both planetary and stellar companions at a very large scale. Finally, this dataset and project represent a first taste of what we can expect for Gaia DR4 -- in terms of sensitivity, challenges around astrophysical false positives, and the massive computational resources needed for this new era of exoplanet population studies.

\section*{Afternoon}
\addcontentsline{toc}{section}{Wednesday -- Afternoon}

\subsection*{A-1502}
\textbf{Theme:} Compact Objects

\hypertarget{abstract-talk-52}{}\label{abstract-page-talk-52}
\subsubsection*{Evidence for a repeating fast radio burst source embedded in a young supernova remnant}
\addcontentsline{toc}{subsubsection}{Evidence for a repeating fast radio burst source embedded in a young supernova remnant}
\textbf{Speaker:} Ayush Pandhi\\
\textbf{Fields:} Compact Objects $\cdot$ Transients

Fast radio bursts (FRBs) are energetic, micro- to millisecond duration radio transients originating from cosmological distances. However, even two decades after their discovery, the specific origins of FRBs remain elusive. In this talk, I will present 3.2 years of observations with the Canadian Hydrogen Intensity Mapping Experiment FRB backend (CHIME/FRB) tracing the evolution of the dynamic local environment around repeating FRB 20220529A by measuring propagation effects imprinted upon its radio signals. We find that the integrated electron column in the FRB environment decreases linearly over multiple years, closely matching theoretical predictions for a FRB source embedded within a young ($\sim$years to a few centuries old), expanding supernova remnant. This scenario is also consistent with the scattering timescale variations and depolarization we observe in FRB 20220529A bursts. This FRB source also undergoes a month-long period in which the local plasma density and magnetic field strength precipitously rise before declining to their original state. I will explore physical scenarios that can explain such extreme, transient changes in the FRB environment, for example magnetized plasma overdensities in the supernova remnant, intervening ejecta or wind from a binary companion, or a magnetar ejecta interactions with a supernova remnant.

\hypertarget{abstract-talk-137}{}\label{abstract-page-talk-137}
\subsubsection*{Crystallization, Convection, and Magnetism in White Dwarf Stars}
\addcontentsline{toc}{subsubsection}{Crystallization, Convection, and Magnetism in White Dwarf Stars}
\textbf{Speaker:} Matias Castro-Tapia\\
\textbf{Fields:} Compact Objects $\cdot$ High Energy Astrophysics

White dwarfs (WDs) are the final evolutionary stage of most stars in the Galaxy. During WD cooling, the strongly coupled plasma in their interiors can form a solid lattice. Such a phase transition induces a compositional buoyant instability in the remaining liquid coupled plasma due to differences in the equilibrium mixture between the two phases. This instability drives convection that mixes the material and can potentially power a dynamo, generating a magnetic field. We use mixing-length theory, WD evolutionary models, and hydrodynamical simulations to investigate crystallization-driven convection in WDs. We find that convection must transition between two regimes of transport efficiency in carbon-oxygen WDs: an initial stage of fast motion and a second stage that resembles the physics of thermohaline mixing. During the efficient regime, we find that a dynamo can be driven for a short time and, if it operates, can generate a magnetic field of megagauss strength embedded in the WD interior. We study the evolution of such magnetic fields and compare our results with observations of magnetic WDs. We find that the crystallization-driven dynamo can only explain WDs with magnetic fields below 4 megagauss for masses of 0.5 to 1 M\_sun. For ultramassive WDs ($>$1.1 M\_sun), we explore oxygen-neon cores with sodium and magnesium impurities, since carbon burning in massive WDs leaves higher traces of these elements. We find that accounting for at least a third species in the phase transition modifies convective properties during crystallization. Consequently, the efficient convection regime lasts one to two orders of magnitude longer, and the mixing zone can extend beyond 80\% of the WD mass coordinate, which would favor a crystallization-driven dynamo in magnetic ultramassive WDs.

\hypertarget{abstract-talk-307}{}\label{abstract-page-talk-307}
\subsubsection*{Cosmic Correlations: How Spins, Masses, and Redshifts Shape the Gravitational Wave Universe}
\addcontentsline{toc}{subsubsection}{Cosmic Correlations: How Spins, Masses, and Redshifts Shape the Gravitational Wave Universe}
\textbf{Speaker:} Chaitanya Khamar\\
\textbf{Fields:} Compact Objects $\cdot$ Gravitational waves $\cdot$ High Energy Astrophysics

Recent gravitational wave observations have tightly constrained the mass and redshift distribution of compact binaries, yet the role of spin within this framework remains uncertain. In particular, it is unclear whether the apparent evolution in spin magnitudes arises primarily from correlations with mass, redshift, or both. In this work, we employ various population models to disentangle these dependencies using O4 gravitational-wave data. Building on the transition mass model, we introduce an analogous transition redshift to test whether spin evolution is more strongly governed by source mass or cosmic epoch. This approach allows us to determine whether high mass, low redshift systems exhibit systematically different spin properties than their low mass counterparts. We further explore alternative parameterizations, including smooth redshift dependent spin evolution within mass bins. From this analysis, we find that the effective spin of a binary black hole is largely dependent on the primary mass and mass ratio. By connecting these population level trends to astrophysical formation channels and cosmic star formation histories, our results provide a window into how black hole spins have evolved over cosmic time and offer predictive insights into the future landscape of gravitational wave populations.

\hypertarget{abstract-talk-262}{}\label{abstract-page-talk-262}
\subsubsection*{Hierarchical mergers may explain correlations in the population of binary black hole mergers}
\addcontentsline{toc}{subsubsection}{Hierarchical mergers may explain correlations in the population of binary black hole mergers}
\textbf{Speaker:} Aditya Vijaykumar\\
\textbf{Fields:} Compact Objects $\cdot$ Gravitational waves

Regardless of their initial spins, the merger of two roughly equal mass black holes (BHs) produces a remnant BH of dimensionless spin 0.69. Such remnants can merge with other BHs in dense stellar environments and produce hierarchical mergers. We find that the latest data from the LIGO-Virgo-KAGRA detectors supports the existence of a subpopulation of such mergers. We characterize the properties of this subpopulation, finding that these mergers have (i) unequal mass ratios, and (ii) relatively larger abundance at high redshift. Our results may explain the correlations of BH spins with mass ratio and redshift found in other works, and have implications for the assembly of star clusters in the Universe.

\subsection*{A-4502}
\textbf{Theme:} Exoplanets - II

\hypertarget{abstract-talk-31}{}\label{abstract-page-talk-31}
\subsubsection*{Dynamical Mass Constraints on Transition Disk Perturbers with Absolute Astrometry}
\addcontentsline{toc}{subsubsection}{Dynamical Mass Constraints on Transition Disk Perturbers with Absolute Astrometry}
\textbf{Speaker:} Dori Blakely\\
\textbf{Field:} Exoplanets

Connecting observed planets with the substructures seen in protoplanetary disks is essential for constraining the formation, evolution, and migration of planetary systems. However, only a handful of (proto)planets have been directly detected around young stars hosting protoplanetary disksso far exclusively through direct imaging and transit photometry. In this talk, I will present the first comprehensive dynamical mass constraints on companions in protoplanetary disk systems using a novel combination of calibrated Gaia and Hipparcos absolute astrometry data. I will present the clear astrometric detection of the known low-mass stellar companion HD 142527 B; newly discovered companions to AB Aurigae and MWC 758; candidate companions to HD 97048, UX Tau A, HD 100546 and CQ Tau; and companion mass upper limits for the remaining targets in our sample. Based on these results, I will show how our dynamical mass constraints can be used to directly connect these inferred companions to the protoplanetary disk substructures seen for several systems. I will also present evidence of excess astrometric noise that affects some protoplanetary disk systems, based on the comparison of models using different subsets of the absolute astrometry data and different noise assumptions. Finally, I will discuss the observability of these newly detected companions using facilities such as VLTI/GRAVITY, and outline the prospects for comprehensive protoplanet demographics that will be possible with Gaia data release 4 (GDR4), based on detailed simulations of GDR4 epoch astrometry of the known protoplanet-hosting systems PDS 70 and WISPIT 2.

\hypertarget{abstract-talk-8}{}\label{abstract-page-talk-8}
\subsubsection*{Searching for Planetary-Mass Objects in Nearby Young Moving Groups using SPHEREx and Citizen Science}
\addcontentsline{toc}{subsubsection}{Searching for Planetary-Mass Objects in Nearby Young Moving Groups using SPHEREx and Citizen Science}
\textbf{Speaker:} Azul Ruiz Diaz\\
\textbf{Field:} Exoplanets

The identification of isolated planetary-mass objects (planemos) is critically needed, as they provide gas giant atmosphere analogs highly suitable to spectroscopic characterization without the glare of a bright host star. Identifying nearby planemos requires locating substellar members of nearby young moving groups; these sparse sub-associations of coeval stars have a common motion through space, and serve as age-dating benchmarks that allows to lift the mass degeneracy in substellar objects. I will present our most recent Bayesian kinematic membership analysis of new substellar objects discovered by citizen scientists through the Backyard Worlds: Planet IX project, supplemented by space-based infrared SPHEREx spectroscopy to identify, validate, and characterize planetary-mass objects in our immediate neighborhood. Our approach leverages the latest version of the BANYAN $\Sigma$ Bayesian classification tool, which includes thousands of young associations. We have so far identified 380 age-calibrated planemos and brown dwarfs with masses $\sim$2 to 88 Mjup, distances $\sim$ 11 to 81 pc, and ages $\sim$10 to 800\,Myr. Of these objects, 80 already have SPHEREx spectra which have allowed us to confirm them as substellar objects. Our results highlight the power of combining citizen science and SPHEREx spectroscopy to uncover and characterize large populations of planemos and brown dwarfs in nearby young moving groups.

\hypertarget{abstract-talk-2}{}\label{abstract-page-talk-2}
\subsubsection*{Unveiling the Nature of Super-Puffs: Early Results from a Panchromatic JWST Atmospheric Spectroscopy Survey}
\addcontentsline{toc}{subsubsection}{Unveiling the Nature of Super-Puffs: Early Results from a Panchromatic JWST Atmospheric Spectroscopy Survey}
\textbf{Speaker:} Michael Radica\\
\textbf{Field:} Exoplanets

With Neptune-like masses but Saturn-like radii, super-puff exoplanets present a fantastic opportunity for spectroscopic atmospheric characterization. With their abnormally large radii and incredibly low densities these planets have been likened to cotton candy, and present significant challenges to conventional models of planet formation. To unravel the origins of super-puffs, we are carrying out a 95 hr survey of the four most observable super-puff planets using JWST's NIRISS, NIRSpec, and MIRI instruments to construct panchromatic optical-to-MIR atmospheric spectra of these objects. In this talk, I will give an overview of some initial results from this survey. The planets in our sample span a range of nearly 1000K in equilibrium temperature allowing us to study disequilibrium chemistry, cloud formation, and limb asymmetry at incredible signal-to-noise over a number of key temperature transitions. These spectra are some of the most precise that JWST will obtain over the entirety of its lifetime, and we are able to robustly disentangle the impacts of hazes and internal heat on the inflated radii of the planets in this sub-population. We also unlock the full chemical inventory of these planets and obtain unprecedented insights into the nature of these most enigmatic of worlds.

\hypertarget{abstract-talk-126}{}\label{abstract-page-talk-126}
\subsubsection*{The Carbon Key: First Detailed Chemical Inventory of the Steam World Planet GJ 9827 d}
\addcontentsline{toc}{subsubsection}{The Carbon Key: First Detailed Chemical Inventory of the Steam World Planet GJ 9827 d}
\textbf{Speaker:} Caroline Piaulet-Ghorayeb\\
\textbf{Field:} Exoplanets

In recent years, JWSTs sensitivity has opened a small sub-Neptune opportunity, paving the way for atmospheric reconnaissance on rocky exo-Earths. The atmospheres of the smallest sub-Neptunes are ideally puffy for observational characterization, but thin enough to carry the imprints of both chemical exchanges with their interiors and atmospheric escape over Gyr timescales. While sub-Neptune mantles may remain molten, resulting in chemical exchanges between the interior and the atmosphere, carbon species are much less soluble than water in the silicate melt, making the atmospheric abundances of carbon species more reliable tracers of planetary elemental makeup relative to water.

The 2-Earth-radii warm (680 K) planet GJ 9827d is the smallest sub-Neptune to be chemically characterized with JWST to date, with NIRISS/SOSS revealing its one-of-a-kind steam world high-metallicity, water-rich atmosphere. I will present results from two new JWST NIRSpec G395H transit observations of GJ 9827 d that provide the first constraints on its carbon budget. The relative abundances of CH4 and CO2 relative to water surprisingly still suggest a water-rich atmosphere, with a deficit in atmospheric carbon compared to expectations from magma ocean-atmosphere interactions. I will discuss this finding in the context of our chemical constraints on the stellar composition, modeling of the structure and chemistry of GJ 9827 ds interior, as well as the possibility of graphite precipitation, early carbon loss, or mantle solidification. The carbon-poor chemistry of GJ 9827 d is at odds with other well-characterized larger sub-Neptunes, suggesting a different formation pathway in line with its unique metal-rich atmosphere.

\hypertarget{abstract-talk-23}{}\label{abstract-page-talk-23}
\subsubsection*{The Fate of Ice Sheets on Chaotically Rotating Exoplanets}
\addcontentsline{toc}{subsubsection}{The Fate of Ice Sheets on Chaotically Rotating Exoplanets}
\textbf{Speaker:} Sarah Silverman\\
\textbf{Field:} Exoplanets

In our search for life beyond the Solar System, the most abundant stars in the Galaxyred dwarfsare the most promising place to look. Most Earth-like exoplanets orbit within the Habitable Zone of these small, dim stars where conditions allow for liquid surface water. The habitability of these planets is threatened by water being lost to space, trapped in their interiors, and frozen in ice sheets. The fate of ice sheets in particular is critical to understand, as ice sheets not only govern the availability of liquid surface water but also interact and evolve with the climate system and the solid planet. Ice sheets have been modeled on synchronously rotating planets in 1:1 spin-orbit resonance. However, gravitational interactions between neighboring planets may prevent synchronous rotation and cause the substellar point to slowly oscillate or circulate across the planet on years-long timescales. Such slow rotation would bring an ice sheet from the nightside into daylight, with global repercussions.

We use a 3D global climate model to simulate planetary climate for near-synchronously rotating planets, then account for key geophysical processes, including basal melting and bedrock deformation, to assess ice sheet stability and evolution. We determine the spatial distribution and thickness of ice sheets and evaluate their impact on water-trapping. This work advances our understanding of the diverse environments that can potentially sustain life beyond Earth and provides insight into how ice sheets, climate, and the solid surface evolve under varied conditions.

\hypertarget{abstract-talk-167}{}\label{abstract-page-talk-167}
\subsubsection*{Seven needles in a haystack: Uncovering the planetary RV signal of the TRAPPIST-1 system with SPIRou and NIRPS}
\addcontentsline{toc}{subsubsection}{Seven needles in a haystack: Uncovering the planetary RV signal of the TRAPPIST-1 system with SPIRou and NIRPS}
\textbf{Speaker:} Alexandrine L'Heureux\\
\textbf{Field:} Exoplanets

The TRAPPIST-1 system is well-known for its seven transiting Earth-sized exoplanets. With ordered and tightly packed planetary orbits around a small M8-type star, the system offers an ideal environment for studying terrestrial planet formation and conducting in-depth atmospheric analyses. These studies rely heavily on the physical parameters of the planets, which have so far been derived primarily from photometric data, with transit timing variations (TTV) used to estimate their masses. As part of the SPIRou Legacy Survey and the NIRPS-GTO, we undertook a radial velocity (RV) follow-up program to directly measure the masses of the TRAPPIST-1 planets. For the first time, we successfully detect the combined planetary RV signal from TRAPPIST-1, confirming the masses inferred from TTV and placing meaningful constraints on the presence of additional non-transiting giant planets in the system. Observing such a faint and active star presents several challenges. I will demonstrate how we isolated the planetary signal from the stellar variability and observing systematics by combining gaussian processes (GP) with the state-of-the-art LBL precision RV framework.

\subsection*{A-3561}
\textbf{Theme:} Education and Public Outreach - I

\hypertarget{abstract-talk-310}{}\label{abstract-page-talk-310}
\subsubsection*{Working Towards a More Inclusive and Rewarding Research Group Culture in Astronomy}
\addcontentsline{toc}{subsubsection}{Working Towards a More Inclusive and Rewarding Research Group Culture in Astronomy}
\textbf{Speaker:} Carolina Cruz-Vinaccia\\
\textbf{Field:} Equity, Diversity and Inclusion

Despite increased attention to equity, diversity, and inclusion (EDI) in Canadian astronomy and physics, astronomers from marginalized groups continue to encounter systemic barriers that affect belonging and persistence in the field. Research groups are central sites of mentorship, training, and professional socialization, making them critical to cultivating more inclusive disciplinary cultures. In this presentation, we describe the development and implementation of an EDI in Action Plan for the Canada Excellence Research Chair (CERC) in Transient Astrophysics research group. The plan embeds equity considerations into the structures and practices that shape group life, with the goal of fostering an inclusive research environment and supporting equitable recruitment, mentorship, and retention of highly qualified personnel. The action plan was developed collaboratively with EDI scholars and informed by relevant social science literature that takes a critical intersectional approach to equity. The plans action areas are aimed at shaping research group culture and supporting the success of all group members through shared values and expectations, supporting community integration, continuing EDI education, strengthening mentorship practices, and establishing mechanisms for accountability. We will highlight three key tools that we have developed that operationalize these goals by establishing transparent norms around mentorship, communication, and participation: 1) Shared Values, 2) Practices \& Expectations, and 3) Onboarding. By sharing both the development process and the outputs, we aim to provide a practical framework that other astronomy research groups can adapt when designing more inclusive and supportive research environments.

\hypertarget{abstract-talk-135}{}\label{abstract-page-talk-135}
\subsubsection*{One Attempt at Building an Inclusive \& Accessible Hybrid Astronomy Conference: FRB 2025}
\addcontentsline{toc}{subsubsection}{One Attempt at Building an Inclusive \& Accessible Hybrid Astronomy Conference: FRB 2025}
\textbf{Speaker:} Alice P. Curtin\\
\textbf{Field:} Equity, Diversity and Inclusion

In a given year, there are hundreds of possible conferences and schools for researchers to attend within the field of astronomy. For example, in 2019 alone, there were over 350 astronomical conferences and summer schools. While meetings are essential for interacting with other researchers and establishing collaborations, many remain only accessible to those with substantial travel funding, flexible schedules, or geographical proximity. This introduces barriers that predominantly affect early career researchers (ECRs) and people from under-resourced regions, limiting the growth, diversity, and sustainability of communities. To address these issues, the yearly Fast Radio conference (with over 200 participants) was held as a fully hybrid conference in Montreal in July 2025. It was designed to prioritize inclusivity and accessibility alongside scientific excellence. In this presentation, I will discuss the goals of FRB 2025, the implementation of those goals, and the feedback received in a post-conference survey. I will advocate for similar hybrid efforts throughout Canada, and discuss how some of the lessons learned could be applied to the virtual CASCA 2027 meeting.

\hypertarget{abstract-talk-3}{}\label{abstract-page-talk-3}
\subsubsection*{A Comparison of Four Treatments to Teach the Seasons}
\addcontentsline{toc}{subsubsection}{A Comparison of Four Treatments to Teach the Seasons}
\textbf{Speaker:} Pierre Chastenay\\
\textbf{Field:} Education and Public Outreach

Refutation texts are texts that begin with a common, explicit misconception, followed by a refutation cue and an explanation based on the currently accepted scientific view (Tippett, 2004). These texts have been shown to promote conceptual change in several areas of science education, including astronomy (Broughton et al., 2010; Mason et al., 2017). However, refutation texts about the seasons have had limited success, partly because the seasons have proven to be one of the most difficult topics to teach and learn in astronomy (Plummer, 2012). We believe that the highly dynamic and 3D nature of this phenomenon is not well conveyed by a simple refutation text, even when accompanied by standard textbook-style illustrations. We propose using a refutation video, showing a dynamic, allocentric view (i.e., the view from space) on the Earth-Sun system and addressing the two most common misconceptions about the seasons (i.e., the variation in the Earth-Sun distance and the Earth's tilt bringing one hemisphere closer to the Sun than the other), which could prove more effective in promoting conceptual change. To test this idea, we used a 2x2 factorial design to compare four different methods of teaching the seasons: a refutation text, an explanatory text (i.e., a standard text from a school textbook), a refutation video, and an explanatory video. Both videos used the texts (refutation and explanatory) as narration. Two hundred participants were randomly assigned to one of the four treatments and were asked to answer seven multiple-choice questions about the seasons before and after the treatment. In this talk, we will present the results of this research, discuss its implication on astronomy teaching, and propose future developments.

\programdayheader{Thursday}
\daypalettebanner{Thursday}

\section*{Morning}
\addcontentsline{toc}{section}{Thursday -- Morning}

\subsection*{A-1502}
\textbf{Theme:} Astrostatistics, Astroinformatics, and AI / Machine Learning

\hypertarget{abstract-talk-251}{}\label{abstract-page-talk-251}
\subsubsection*{Beyond the Power Spectrum: A New Framework for Non-Stationary Fields with Applications to Light-Cone Effects in Line Intensity Mapping}
\addcontentsline{toc}{subsubsection}{Beyond the Power Spectrum: A New Framework for Non-Stationary Fields with Applications to Light-Cone Effects in Line Intensity Mapping}
\textbf{Speaker:} Mattéo Blamart\\
\textbf{Fields:} Astrostatistics, Astroinformatics, and AI / Machine Learning $\cdot$ Cosmology and Early Universe

Modern cosmological surveys cover extremely large volumes and map fluctuations on scales reaching gigaparsecs. As a result, it is longer a valid assumption to ignore cosmological evolution along the line of sight from one end of the survey to the other one. When extracting cosmological information, the power spectrum becomes suboptimal because it relies on the assumption of translationnal invariance of the observed field. For example, during the Epoch of Reionization (EoR), the 21cm brightness temperature field of near its end exhibits statistical properties that differ significantly from those at its beginning. To overcome these limitations, we have developed a novel non-Fourier basis. Our work demonstrates that using this new basis integrated in a new summary statistics yields tighter constraints on astrophysical parameters compared to the traditional power spectrum. Additionally, we provide an illustrative example and a practical guide for applying this basis in the context of a realistic forecast for interferometers such as the Hydrogen Epoch of Reionization Array (HERA).

\hypertarget{abstract-talk-xx2}{}\label{abstract-page-talk-xx2}
\subsubsection*{Emulating Black Hole-Neutron Star Merger Disk Ejecta}
\addcontentsline{toc}{subsubsection}{Emulating Black Hole-Neutron Star Merger Disk Ejecta}
\textbf{Speaker:} Joe Bardwell\\
\textbf{Fields:} Astrostatistics, Astroinformatics, and AI / Machine Learning $\cdot$ Gravitational waves and Transients

Neutron star mergers (NSMs) are multimessenger sources and the only confirmed production site of r-process elements. A central challenge in interpreting their signals is modeling the ejecta that powers the kilonova transient. The post-merger accretion disk is a major source of this material, but ejecta production is complex and requires long simulation timescales. Here I will discuss the development of a machine learning emulator of disk ejecta, designed for rapid parameter inference from future multimessenger observations. By providing a high-quality, multi-variable interpolator, this tool offers significant improvement over simpler prescriptions currently in use. As a first step, we focus on black hole-neutron star mergers. These systems generate accretion disks that contain the key physical ingredients over a well-defined parameter space, while leaving out complexities such as the hypermassive neutron star phase. The resulting neural network framework, based on a large training set of viscous hydrodynamic simulations, predicts key scalar output quantities with error margins of a few percent, while dramatically reducing computational costs: a single function evaluation is completed in under 0.2 core-seconds, compared to approximately 3 core-years for a full simulation. Our results provide a scalable framework for future extensions into the more complex NSM regime.

\hypertarget{abstract-talk-227}{}\label{abstract-page-talk-227}
\subsubsection*{Revealing the intrinsic population of fast radio burst observables through injection analysis}
\addcontentsline{toc}{subsubsection}{Revealing the intrinsic population of fast radio burst observables through injection analysis}
\textbf{Speaker:} Kyle McGregor\\
\textbf{Fields:} Astrostatistics, Astroinformatics, and AI / Machine Learning $\cdot$ Transients

The release of a second catalog of fast radio bursts (FRBs) by the CHIME/FRB collaboration marks a major advance in the statistical study of these enigmatic cosmic transients. As with any astrophysical survey, however, the observed population is shaped by the instruments detection biases. A quantitative characterization of these selection effects is therefore imperative for interpreting the intrinsic FRB population. We probe CHIME/FRBs selection function by injecting a large, synthetic FRB population with known properties into the telescopes live data stream, allowing us to model the probability of detection and non-detection across the multidimensional FRB parameter space. Using these injections, we construct a parametric statistical model of the survey selection function and apply it to correct observed FRB distributions in the recent CHIME/FRB catalog 2 data release, which provides nearly an order of magnitude more events than previous FRB catalogs. I will present the methodology used to infer this selection function, discuss the resulting bias of CHIME/FRB to different burst properties, and show how these corrections enable improved constraints on the underlying intrinsic distributions from which the observed FRB population derives. In particular, this analysis indicates that CHIME/FRB is likely missing a substantial population of bursts with scattering timescales above 10 ms, highlighting the importance of accounting for selection biases when drawing astrophysical conclusions from FRB surveys.

\hypertarget{abstract-talk-68}{}\label{abstract-page-talk-68}
\subsubsection*{Methane from the Mantle: Geophysical Pathways to Prebiotic Atmospheres}
\addcontentsline{toc}{subsubsection}{Methane from the Mantle: Geophysical Pathways to Prebiotic Atmospheres}
\textbf{Speaker:} Nelles Henderson\\
\textbf{Field:} Exoplanets

A key question in the search for life elsewhere in the universe is what planetary conditions are necessary to form life. With the groundbreaking James Webb Space Telescope, we can now observe exoplanet atmospheres directly, opening the possibility of detecting an exoplanet with all the necessary building blocks for life. A methane-rich, reducing atmosphere is likely necessary to generate high abundances of key prebiotic molecules on the surface of a terrestrial planet. Such worlds are typically not massive enough to retain their primary atmospheres, so the compositions of their secondary atmospheres are determined to first order by mantle outgassing. I will present new results obtained using the PROTEUS framework to simulate the coupled evolution of the atmospheres and interiors of rocky planets. PROTEUS is a self-consistent 1D numerical model that accounts for interior cooling, volatile exchange between the mantle and atmosphere, equilibrium atmospheric chemistry, and loss to space. I will share results from simulations spanning one billion years of evolution, from the onset of magma ocean crystallization to long-term climate equilibrium. I will present our analysis of the effect of mantle redox state on atmospheric composition, with a particular focus on the geophysical pathways to methane production. I will also discuss how bulk elemental abundance ratios set by planet formation further control atmospheric chemistry. This detailed treatment of interior-atmosphere coupling on rocky planets ultimately links planetary interiors and formation theory with the detections of exoplanet atmospheres made in state-of-the-art observations today.

\hypertarget{abstract-talk-330}{}\label{abstract-page-talk-330}
\subsubsection*{Dynamical SMBH Mass Measurement in Obscured Galaxies via Joint SMBH-Stellar Mass Modeling}
\addcontentsline{toc}{subsubsection}{Dynamical SMBH Mass Measurement in Obscured Galaxies via Joint SMBH-Stellar Mass Modeling}
\textbf{Speaker:} Michael Yantovski-Barth\\
\textbf{Fields:} Astrostatistics, Astroinformatics, and AI / Machine Learning $\cdot$ Cosmology and Early Universe $\cdot$ Galaxy $\cdot$ Supermassive black hole

Over the past decade, ALMA has revolutionized dynamical measurements of supermassive black hole (SMBH) masses in nearby galaxies, providing some of the strongest constraints on SMBH-galaxy co-evolution. In principle, the magnification due to strong gravitational lensing enables similar SMBH mass measurements out to redshifts of z$\sim$2-7, thereby probing SMBH growth in the early universe. However, modeling the dynamics of lensed galaxies is often hindered by strong dust obscuration and bright active galactic nuclei (AGN), which prevent direct observation of the distributions of the stars. Without stellar-light information, current dynamical models struggle to disentangle the gravitational signatures of the stars and the SMBHs. To address this, we introduce a generalized mathematical framework that models the molecular gas dynamics arising from complex stellar mass distributions and enables them to be fitted directly to kinematical data. We incorporate these stellar mass models into SuperMAGE, a fully differentiable dynamical modeling simulator, and introduce VisCube, a visibility-space gridder that supports uncertainty quantification for radio interferometry. We unite these tools in a Bayesian inference pipeline to measure SMBH masses directly in visibility space while simultaneously constraining the stellar mass distributions even in dust-obscured and AGN-dominated galaxies. We validate our approach by reanalyzing the SMBH mass measurement in the galaxy NGC 4697: our results demonstrate that marginalizing over flexible stellar mass models can mitigate SMBH mass biases induced by AGN light, thereby enabling more accurate comparisons between AGN luminosities and SMBH masses. We conclude by outlining ongoing applications of our methodology to high-redshift lensed galaxies.

\hypertarget{abstract-talk-317}{}\label{abstract-page-talk-317}
\subsubsection*{Joint Source and Lens-Mass Inference in Strong Gravitational Lensing with Diffusion Models}
\addcontentsline{toc}{subsubsection}{Joint Source and Lens-Mass Inference in Strong Gravitational Lensing with Diffusion Models}
\textbf{Speaker:} Gabriel Missael Barco\\
\textbf{Fields:} Astrostatistics, Astroinformatics, and AI / Machine Learning $\cdot$ Galaxy $\cdot$ Gravitational waves

Strong gravitational lensing is a powerful probe of galaxy mass distributions and a unique window onto distant sources, but reconstructing the unlensed source from strong lens observations is a challenging and ill-posed inverse problem. Diffusion models, a class of generative machine learning models, provide expressive, data-driven priors for source reconstruction. Initial diffusion-based approaches assume that the lens mass distribution, and therefore the forward operator, is known. We relax this assumption by jointly inferring the unlensed source and a parametric lens-mass profile from the lensed image. This places source reconstruction with diffusion-model priors and lens modelling in a unified probabilistic framework. In simulated observations, the resulting reconstructions yield residuals consistent with the observational noise, and the marginal posteriors of the lens parameters recover the true values without detectable systematic bias. These results show that joint inference with diffusion-model priors can support accurate source reconstruction and lens-parameter inference. More broadly, this work is a step toward practical Bayesian pipelines for strong-lensing analysis of real data from next-generation surveys with expressive, data-driven priors.



\hypertarget{abstract-talk-35}{}\label{abstract-page-talk-35}
\subsubsection*{Robust Simulation-Based Inference for the Interstellar Medium Using Generative AI}
\addcontentsline{toc}{subsubsection}{Robust Simulation-Based Inference for the Interstellar Medium Using Generative AI}
\textbf{Speaker:} Duo Xu\\
\textbf{Fields:} Astrostatistics, Astroinformatics, and AI / Machine Learning $\cdot$ Gas, dust and the ISM

Simulation-Based Inference (SBI) is rapidly becoming a cornerstone of modern astrophysics, yet traditional discriminative machine learning models frequently struggle with physical interpretability and out-of-distribution (OOD) generalization when applied to real observational data. We present a novel, robust framework leveraging generative artificial intelligencespecifically Denoising Diffusion Probabilistic Models (DDPM) and the Diffusion Schrdinger Bridge (DSB)to overcome these limitations in the study of the interstellar medium (ISM). Unlike pixel-to-pixel discriminative models, our paired Astro-DSB approach and multi-channel DDPM architectures are explicitly designed to align the learned statistical distributions with the underlying dynamical systems. We demonstrate the efficacy of this generative approach across several highly complex, non-linear inverse problems. First, we show how a 3-channel DDPM can accurately infer magnetic field strengths by fusing synthetic observables (column density, polarization angles, and LOS velocity dispersion), outperforming classical DavisChandrasekharFermi methods and maintaining strict robustness on unseen OOD simulation data. Second, we highlight new extensions of these diffusion models to infer challenging 3D intrinsic properties, including volume density and the interstellar radiation field (ISRF), directly from sparse 2D projections. By evaluating the learning process and predictive performance on both physical simulations and real observations, we demonstrate that generative AI offers a highly reliable, physics-aware pathway to bridge the gap between complex simulations and astronomical data.

\hypertarget{abstract-talk-xx}{}\label{abstract-page-talk-xx}
\subsubsection*{DOROTHY: Stellar Catalogue for 13 Million Milky Way Stars Derived Using a Machine Learning Pipeline}
\addcontentsline{toc}{subsubsection}{DOROTHY: Stellar Catalogue for 13 Million Milky Way Stars Derived Using a Machine Learning Pipeline}
\textbf{Speaker:} Diego Martinez Noriega and Isabelle Laing\\
\textbf{Fields:} Astrostatistics, Astroinformatics, and AI / Machine Learning $\cdot$ Stars and Stellar cluster formation 

Large spectroscopic surveys now provide spectra for millions of stars, enabling data-driven approaches to infer stellar properties at an unprecedented scale. However, most deep-learning models are restricted to single-survey applications due to differences in spectral resolution and wavelength coverage. We present DOROTHY, a unified deep neural network that predicts 11 stellar parameters and chemical abundances, including effective temperature, surface gravity, and metallicity, together with their associated uncertainties, from spectra spanning four spectroscopic surveys: DESI DR2, BOSS DR19, and LAMOST DR11 (low- and medium-resolution). The network is trained on a dataset of ? 240,000 stars with spectra in at least one of the four surveys, and measurements from either APOGEE DR17 or GALAH DR4, allowing us to leverage their high-resolution labels to derive parameters from lower-resolution spectra. We demonstrate that DOROTHY achieves high predictive accuracy, with mean absolute residuals of <0.05 dex for most abundances. We release a unified catalogue of predicted stellar parameters and chemical abundances, their associated uncertainties, and four quality flags for ~ 13.6 million stars across multiple major surveys. The DOROTHY framework unifies data analysis across spectroscopic surveys, enabling larger-scale studies of the chemical evolution of the Milky Way.

\subsection*{A-4502}
\textbf{Theme:} Next generation observatories - II

\hypertarget{abstract-talk-49}{}\label{abstract-page-talk-49}
\subsubsection*{Measuring gravitational lensing time delays with quantum information processing}
\addcontentsline{toc}{subsubsection}{Measuring gravitational lensing time delays with quantum information processing}
\textbf{Speaker:} Shiming Gu\\
\textbf{Fields:} Exoplanets $\cdot$ Instrumentation $\cdot$ Transients

The gravitational fields of astrophysical bodies bend the light around them, creating multiple paths along which light from a distant source can arrive at Earth. Measuring the difference in photon arrival time along these different paths provides a means of determining the mass of the lensing system, which is otherwise difficult to constrain. This is particularly challenging in the case of microlensing, where the images produced by lensing cannot be individually resolved; existing proposals for detecting time delays in microlensed systems are significantly constrained due to the need for large photon flux and the loss of signal coherence when the angular diameter of the light source becomes too large. In this work, we propose a novel approach to measuring astrophysical time delays. Our method uses exponentially fewer photons than previous schemes, enabling observations that would otherwise be impossible. Our approach, which combines a quantum-inspired algorithm and quantum information processing technologies, saturates a provable lower bound on the number of photons required to find the time delay. Our scheme has multiple applications: we explore its use both in calibrating optical interferometric telescopes and in making direct mass measurements of ongoing microlensing events. To demonstrate the latter, we present a fiducial example of microlensed stellar flare sources in the Galactic Bulge. Though the number of photons produced by such events is small, we show that our photon-efficient scheme opens the possibility of directly measuring microlensing time delays using existing and near-future ground-based telescopes.

\hypertarget{abstract-talk-265}{}\label{abstract-page-talk-265}
\subsubsection*{The Colibri Telescope Array: A Robotic Facility for Time-Domain Astronomy}
\addcontentsline{toc}{subsubsection}{The Colibri Telescope Array: A Robotic Facility for Time-Domain Astronomy}
\textbf{Speaker:} Toni Cordeiro de Almeida\\
\textbf{Field:} Instrumentation

The Colibri Telescope Array (Colibri) is a robotic facility designed to probe the time-domain Universe. It consists of three 0.5-m f/3 telescopes manufactured by Hercules Telescopes in Montréal, each equipped with a high-cadence (up to 40 fps) FLI Kepler KL4040 sCMOS camera, providing a field of view of 1.43 1.43. Located at Elginfield Observatory in Ontario, Canada, Colibri is designed to detect kilometer-sized trans-Neptunian objects through serendipitous stellar occultations, where rapid imaging and coincidence multi-site detection are necessary to distinguish false positives from real detection. However, Colibri is also capable of longer-exposure astronomical observations, with guiding soon to be available on two of the three telescopes. In addition to its optical capabilities, Colibri is expanding into the short-wavelength infrared (900 - 1700 nm) with the integration of a Raptor Photonics Owl 1280 InGaAs camera. Robotic observations are made possible through the use of commercial off-the-shelf hardware and software communicating via the ASCOM protocol. We present the Colibri instrumentation and demonstrate its current capabilities. The Colibri Telescope Array team invites the Canadian astronomical community to collaborate and to submit requests for telescope time.

\hypertarget{abstract-talk-156}{}\label{abstract-page-talk-156}
\subsubsection*{CCAT: The Fred Young Submillimeter Telescope (FYST) nears first light on Cerro Chajnantor}
\addcontentsline{toc}{subsubsection}{CCAT: The Fred Young Submillimeter Telescope (FYST) nears first light on Cerro Chajnantor}
\textbf{Speaker:} Douglas Scott\\
\textbf{Fields:} Instrumentation $\cdot$ Next generation observatories

The Fred Young Submillimeter Telescope (FYST) is a 6-m aperture wide-field-of-view telescope situated at 5600 m elevation near the summit of Cerro Chajnantor in northern Chile, which is scheduled for completion in July of this year. It is run by the CCAT Consortium, which includes a significant component of Canadian astronomers, funded through the CFI and the Provinces of Ontario, BC and Alberta. The telescope's properties, its site, and its state-of-the-art instrumentation promise unrivalled mapping speed for (sub)mm-wave surface brightness studies over large fields. The science goals include: tracing the formation and evolution of star-forming galaxies from the epoch of reionization to the cosmic peak of star-formation activity through wide-field dust continuum surveys and broad-band [CII] line imaging; constraining thermodynamics and feedback in galaxy clusters through the Sunyaev-Zeldovich effects; improving constraints on primordial gravitational waves through precision removal of polarized foregrounds; and tracing local star-formation processes through velocity-resolved spectroscopy over a large range of scales in the Galaxy. FYST will be sequentially receiving its first-light instrumentation, a submm heterodyne receiver (MiniCHAI) and a wide-field submm continuum receiver (Prime-Cam) for first light in mid-2026. This talk will include a brief review of the overall science objectives addressed by CCAT, the current state of the instrumentation and telescope, and finish with a discussion of the first-light and early science programs that are planned.

\hypertarget{abstract-talk-290}{}\label{abstract-page-talk-290}
\subsubsection*{The Canadian Galactic Emission Mapper: Overview, Optical Design, and On-Sky Characterization}
\addcontentsline{toc}{subsubsection}{The Canadian Galactic Emission Mapper: Overview, Optical Design, and On-Sky Characterization}
\textbf{Speaker:} Joshua MacEachern\\
\textbf{Fields:} Cosmology and Early Universe $\cdot$ Gas, dust and the ISM $\cdot$ Instrumentation

Gravitational waves from cosmic inflation may have left a detectable signature in the parity-odd, "B-mode" component of the polarization of the Cosmic Microwave Background (CMB). Detecting B-modes in the CMB would be "smoking gun" evidence for inflation and would probe some of the highest-energy physics in the known universe. However, current experiments have placed stringent upper limits on B-modes. If B-modes are present in the CMB, the signal is extremely faint and is dominated by polarized Galactic foregrounds at all frequencies. It is therefore essential to map polarized foregrounds with high precision to enable a detection of CMB B-modes. The Canadian Galactic Emission Mapper (CGEM) is a new 4m, single-dish, radio telescope at the Dominion Radio Astrophysical Observatory that is mapping polarized Galactic synchrotron emission from 8-10 GHz over the Northern sky. CGEM will greatly improve models of polarized CMB foregrounds and will hence be an important aid to current and future B-mode experiments. In this talk, I'll give an overview of CGEM's science goals. I'll then describe how we designed this purpose-built instrument to measure the sky with minimal polarization systematics. I'll highlight in particular the polarization purity of the optical design, and will describe how we've optimized the optics to minimize polarization leakage. I'll then showcase early commissioning observations from a pathfinder version of CGEM. On-sky measurements of the Sun, satellites, and other sources show a stunning agreement with our optical models. Amazingly, we're able to use Earth-observation satellites seen by CGEM to probe the beam to 0.01\% of the maximum. The excellent on-sky performance at this early stage in the experiment shows great promise for future science with CGEM.

\hypertarget{abstract-talk-43}{}\label{abstract-page-talk-43}
\subsubsection*{Systematics-aware forecasts for HI intensity mapping with CHORD}
\addcontentsline{toc}{subsubsection}{Systematics-aware forecasts for HI intensity mapping with CHORD}
\textbf{Speaker:} Sophia Rubens\\
\textbf{Fields:} Cosmology and Early Universe $\cdot$ Instrumentation $\cdot$ Next generation observatories

As fruitful as neutral hydrogen (HI) observations with existing observatories have been, upcoming instruments stand to facilitate even more precise measurements. This opens the door to improved constraints on the dark energy equation of state. A telescope well-positioned to take such measurements in the pre-SKA era is the Canadian Hydrogen Observatory and Radio-transient Detector (CHORD), an under-construction drift-scan radio interferometer planned to observe between 300 and 1500 MHz (z=0-3.7) which will comprise 512 six-metre dishes at the Dominion Radio Astrophysical Observatory (DRAO). When targeting unprecedentedly precise measurements of how HI traces the underlying matter distribution, we must relax typical assumptions of instrumental systematics having a negligible effect on cosmological measurements and forecasts. I will present the results of end-to-end simulations that provide forecasts on CHORD's ability to measure cosmological parameters, including \$\textbackslash{}Omega\_m\$, in the presence of realistic systematics informed by mm-scale dish surface imperfections and feed displacements.

\hypertarget{abstract-talk-110}{}\label{abstract-page-talk-110}
\subsubsection*{The Future of Precision Astronomy: The Best Telescope is All the Telescopes}
\addcontentsline{toc}{subsubsection}{The Future of Precision Astronomy: The Best Telescope is All the Telescopes}
\textbf{Speaker:} Kevin McKinnon\\
\textbf{Fields:} Astrostatistics, Astroinformatics, and AI / Machine Learning $\cdot$ Next generation observatories

The next generation of observatories -- including Euclid, Roman, and Rubin -- will deliver deep astrometric data across the sky, enabling high-precision kinematic studies of the faintest stars in the Milky Way (MW), the Local Group, and beyond. These faint populations, largely inaccessible to Gaia, are crucial for near-field cosmology: they hold the key to resolving the MW's stellar halo in high resolution, reconstructing the Local Group's dynamical history, discovering and characterizing distant stellar streams, and revealing the nature of dark matter. However, the best possible astrometry will come from harnessing the strengths of all telescopes simultaneously. In this talk, I will present ongoing work combining Hubble and James Webb images with Gaia data to measure significantly improved positions, parallaxes, and proper motions (PMs), effectively unlocking next-generation precision a decade early. Tests with nearby dwarf galaxies yield PMs up to 50 times more precise than Gaia alone; in some cases, we can extend our high-quality PMs to stars hundreds of times fainter than Gaia's limit. I will highlight exciting applications of this approach within the Local Group, focusing on the distant MW stellar halo, stellar streams, and MW/M31 satellite kinematics. I will discuss generalizations of our techniques to ground-based archival data -- such as legacy and contemporary CFHT surveys -- as well as next-generation telescopes. Finally, I will conclude with simulation-driven predictions about the future of combined-telescope astrometry and high-precision PMs in the era of Rubin, Euclid, and Roman.

\hypertarget{abstract-talk-157}{}\label{abstract-page-talk-157}
\subsubsection*{Plans for CFHT's next decade}
\addcontentsline{toc}{subsubsection}{Plans for CFHT's next decade}
\textbf{Speaker:} Nadine Manset\\
\textbf{Field:} Instrumentation

The Canada-France-Hawaii Telescope (CFHT) 10-year strategic plan, aligned with its current lease through 2033, focuses on ensuring operational sustainability while advancing key scientific and community initiatives. Strategic priorities include the development of the new Wenaokeao instrument, the launch of a scientifically ambitious multi-year Community Surveys (CS) program using MegaCam and Wenaokeao, and meaningful engagement with the residents and Hawaiians of the Big Island and the state to help shape a vision for CFHT and astronomy on Maunakea beyond 2033. Targeted deferred maintenance activities, both recent and planned, are strengthening core infrastructure and operational resilience, critical foundations for the successful execution of the Community Surveys. In parallel, ongoing discussions with the Mauna Kea Stewardship and Oversight Authority and members of the local community are helping define a sustainable and collaborative path forward for the post-2033 era. The development of our mutual stewardship model will include the astronomical and local communities.

\hypertarget{abstract-talk-194}{}\label{abstract-page-talk-194}
\subsubsection*{Advancing Coherent Differential Imaging: Laboratory Gains and First On-Sky Results with the SPIDERS Self-Coherent Camera}
\addcontentsline{toc}{subsubsection}{Advancing Coherent Differential Imaging: Laboratory Gains and First On-Sky Results with the SPIDERS Self-Coherent Camera}
\textbf{Speaker:} Kaitlyn Hessel\\
\textbf{Field:} Instrumentation

Growing efforts to directly image and characterize Earth-like exoplanets orbiting nearby stars has accelerated the development of new instruments and techniques aimed to push the current limits of direct imaging and achieve unprecedented levels of contrast in high-contrast imaging. A post-processing technique called Coherent Differential Imaging (CDI) offers a powerful solution by using the incoherence of planet and stellar light to remove the stellar speckle field without suffering from self-subtraction or distortion of planetary signals, giving it theoretically infinite contrast in monochromatic light. Using a Self-Coherent Camera (SCC) integrated into the Subaru Pathfinder Instrument for Detecting Exoplanets and Retrieving Spectra (SPIDERS), we have demonstrated contrast gains of up to 30x in laboratory testing and conducted the first on-sky demonstration of CDI with an SCC at the Subaru Telescope, even under challenging weather conditions. These results validate the technique outside the lab and highlight its future capabilities for upcoming instruments and flagship missions. I will present our recent on-sky results, key lessons learned and the path forward for using CDI in future instruments including the Gemini Planet Imager 2 and NASAs Habitable Worlds Observatory.

\hypertarget{abstract-talk-151}{}\label{abstract-page-talk-151}
\subsubsection*{Instrumentation highlights, R\&D and Canadian collaborations at NRC Herzberg}
\addcontentsline{toc}{subsubsection}{Instrumentation highlights, R\&D and Canadian collaborations at NRC Herzberg}
\textbf{Speaker:} Alan McConnachie\\
\textbf{Field:} Instrumentation

I present the status of some of the of the key facility instrumentation projects underway at NRC Herzberg relevant for Canadian Long Range Plan priorities and Canadian-led collaborations in which Herzberg participates. These include CASTOR technologies for a UVMOS instrument, next generation correlators for SKA-mid and ALMA, and engagement in ELT instrumentation. I also emphasise the R\&D activities and pan-Canadian collaborations that make these contributions possible.

\hypertarget{abstract-talk-172}{}\label{abstract-page-talk-172}
\subsubsection*{OMM : Boosting small telescope performances with EMCCD!}
\addcontentsline{toc}{subsubsection}{OMM : Booting small telescope performances with EMCCD!}
\textbf{Speaker:} Lison Malo\\
\textbf{Field:} Instrumentation

Electron multiplying CCDs (EMCCDs), particularly when operated in photon counting mode, open new possibilities for high cadence and low flux astronomical observations. Despite its modest size, the 1.6m telescope of the Observatoire du Mont-Mégantic (OMM) can achieve performance comparable to, or even surpassing, larger facilities when observations are limited by read noise. This presentation will highlight the growing integration of EMCCD technology at OMM through three complementary instruments. We will first present PESTO, an EMCCD camera in operation for several years, which has enabled fast photometry, the monitoring of variable objects, and high-sensitivity imaging, demonstrating the relevance of EMCCDs for low read noise observations. We will then introduce VROOMM, an optical high-resolution spectrograph currently in development using a 4k4k EMCCD (E2V CCD282) driven by Nuvu CCCP controller, designed to deliver unprecedented high cadence spectroscopic capabilities for a telescope of this class. Finally, we will describe the forthcoming integration of EMCCDs into the SPIOMM imaging Fourier transform spectrometer (aka new THÉSÉE instrument), paving the way for more sensitive hyperspectral observations. We will also discuss challenges associated with high cadence EMCCD instruments, particularly data management and archiving, which require close collaboration with the CADC to adapt pipelines and storage systems. Together, these developments position OMM as a unique testbench for next generation low noise EMCCD instrumentation.

\subsection*{A-3561}
\textbf{Theme:} Supermassive black hole

\hypertarget{abstract-talk-18}{}\label{abstract-page-talk-18}
\subsubsection*{AGN Feedback Kiloparsec-Scale Signatures in Cosmological Simulations}
\addcontentsline{toc}{subsubsection}{AGN Feedback Kiloparsec-Scale Signatures in Cosmological Simulations}
\textbf{Speaker:} Marine Prunier\\
\textbf{Fields:} Galaxy $\cdot$ Galaxy Clusters $\cdot$ Supermassive black hole

AGN feedback plays a central role in shaping the thermodynamic structure of galaxy clusters, driving multiphase outflows that redistribute energy and metals across the intracluster medium (ICM). Capturing these complex processes in cosmological simulations is challenging, particularly the generation of observable small, kiloparsec-scale signatures such as X-ray cavities, shocks, and sound waves. In this talk, I will review recent advances in state-of-the-art cosmological simulations in reproducing AGN-driven structures that impact the hot ICM on kiloparsec scales. I will discuss methodologies for consistently comparing X-ray cavities and weak shocks in simulations with observations. Finally, I will highlight two key implications of these advancements: first, opening a new observational window to study the role of AGN feedback in a full cosmological context; and second, demonstrating that careful comparisons between these small-scale X-ray structures in simulations and observations are essential for validating SMBH subgrid models.

\hypertarget{abstract-talk-5}{}\label{abstract-page-talk-5}
\subsubsection*{Resolving the Hidden Kinematics of Galaxy Cluster Filaments with Ultra-High-Resolution Spectroscopy}
\addcontentsline{toc}{subsubsection}{Resolving the Hidden Kinematics of Galaxy Cluster Filaments with Ultra-High-Resolution Spectroscopy}
\textbf{Speaker:} Benjamin Vigneron\\
\textbf{Fields:} Galaxy $\cdot$ Galaxy Clusters $\cdot$ Supermassive black hole

We present new insights into the kinematics, ionization, and geometry of the optical filamentary nebulae surrounding the brightest cluster galaxies NGC 4696 and NGC 1275, based on extremely high spectral resolution Gemini/GHOST (R = 60,000) spectroscopy and archival CFHT/ESPaDOnS spectro-polarimetric data (R = 68,000). These optical filaments form in galaxy clusters where active galactic nuclei (AGN) feedback suppresses large-scale cooling of the hot ($\sim10^7$ K) X-ray emitting intra-cluster medium (ICM), yet allows the survival of cooler ($\sim10^4$ K) ionized optical gas.

Targeting key regions in the filaments, our GHOST and ESPaDOnS observations resolve the main optical emission lines (H$\alpha$, [O III], [N II], [S II]) into multiple kinematic components for the first time. We detect extremely narrow emission features, with velocity dispersions below 40 km s$^{-1}$, embedded within broader line structures. These components, inaccessible to past lower-resolution studies, reveal a complex superposition of filamentary substructures spanning a wide dynamical range: from bright, turbulent filaments closely coupled to AGN activity to faint, quiescent structures likely shielded from the surrounding ICM.

We combine these kinematic constraints with line-ratio diagnostics to probe the physical conditions of the ionized gas and present a first exploration of the polarimetric properties of optical filaments in the Perseus cluster with ESPaDOnS polarimetric data. Together, these results reveal a clear dichotomy between inner and outer filaments in both kinematics and ionization state, tentatively pointing toward distinct formation origins. Our study thus demonstrates that ultra-high spectral resolution spectroscopy opens a new avenue of studies on the structure of galaxy cluster filaments.

\hypertarget{abstract-talk-184}{}\label{abstract-page-talk-184}
\subsubsection*{New Infrared BPT Diagnostics for JWST: Probing Warm Gas in Massive Galaxies}
\addcontentsline{toc}{subsubsection}{New Infrared BPT Diagnostics for JWST: Probing Warm Gas in Massive Galaxies}
\textbf{Speaker:} Olivia Pereira\\
\textbf{Fields:} Galaxy $\cdot$ Galaxy Clusters $\cdot$ Supermassive black hole

One of the central questions in astrophysics is how galaxies grow and regulate themselves over cosmic time. This process is governed by a complex cycle of gas cooling, star formation, and energy released by central supermassive black holes. The most massive galaxies in the Universe, Brightest Cluster Galaxies (BCGs) at the centers of galaxy clusters, offer a unique window into this cosmic ecosystem, where enormous reservoirs of hot gas, multiphase filaments, and powerful black hole feedback operate over scales from hundreds of parsecs to hundreds of kiloparsecs. In this talk, I will present new JWST/NIRSpec Integral Field Unit (IFU) observations of the central 640 $\times$ 640 pc of NGC 4696, the prototypical BCG in the Centaurus cluster. These data reveal more than 20 atomic and molecular emission lines, including a rich set of H$_2$ rovibrational transitions that trace spatial variations in excitation mechanisms at unprecedented resolution. Using Cloudy simulations, I develop new infrared diagnostics, analogous to optical BPT diagrams, optimized for JWST observations. These results provide new insights into the physical state and excitation of warm gas in massive galaxies and demonstrate the broader potential of JWST spectroscopy for probing feedback, cooling, and star-black hole coupling across different environments. This work highlights Canadas significant role in JWST science and its contribution to advancing infrared galaxy evolution research, fostering pride in our scientific community.

\hypertarget{abstract-talk-329}{}\label{abstract-page-talk-329}
\subsubsection*{A Census of Massive Compact Galaxies from Euclid}
\addcontentsline{toc}{subsubsection}{A Census of Massive Compact Galaxies from Euclid}
\textbf{Speaker:} Ivana Damjanov\\
\textbf{Field:} Galaxy

The Euclid space mission is revolutionizing our view of galaxy evolution, mapping the structural and stellar population properties of unprecedentedly large galaxy samples across cosmic time. In this talk, I will focus on a critical population for testing mass assembly models: massive compact galaxies, systems with the mass of the Milky Way compressed into a region the size of its bulge. Thought to be the ancient, dense cores of today's most massive galaxies, they provide a powerful test of the two-phase formation scenariothe idea that massive galaxies form their compact cores rapidly in the early universe and later assemble their outer envelopes through mergers and accretion. Leveraging Euclid's wide-field, homogeneous imaging, I will present the properties of statistically robust samples of these objects across the broadest redshift range ever studied. Detailed analysis of massive compact systems in Euclid will yield new constraints on their abundance, internal properties, and environmental preferences. These results will strengthen the link between the high-redshift galaxy population and the cores of local massive ellipticals, providing a robust observational anchor for the two-phase formation scenario.

\hypertarget{abstract-talk-203}{}\label{abstract-page-talk-203}
\subsubsection*{statmorph-lsst: how robust are our measurements of galaxy morphology?}
\addcontentsline{toc}{subsubsection}{statmorph-lsst: how robust are our measurements of galaxy morphology?}
\textbf{Speaker:} Elizaveta Sazonova\\
\textbf{Field:} Galaxy

Large upcoming surveys - Rubin, Euclid, SKA, Roman, and more - are beginning to deliver exquisite multi-band imaging of galaxies across cosmic time, letting us trace stellar populations and multi-phase gas in tandem. This imaging can be instrumental in understanding how galaxy structure evolves not just as a whole, but also in different components. However, our measurements of galaxy structure are strongly dependent on image quality, which will vary both across surveys and as a function of source redshift, and so combining these datasets will be a challenge.

Using 189 nearby galaxies observed with HST, we created a dataset of 64,000 augmented images at varying resolutions and depths, and measured their morphology using an adapted morphology library, statmorph-lsst, which is developed as a Rubin contribution from Canada. We systematically characterized how every morphological metric responds to the changes in image quality. I will discuss the biases and common pitfalls in common morphological metrics, the reasons behind them, and potential remedies or corrections. Our results show that some of the puzzling results from JWST analyses of high-redshift galaxies such as an apparent decrease in Sersic and concentration indices in early-type galaxies can be explained purely by observational biases. I will also mention two new metrics introduced in statmorph-lsst. These tools can be applied not only to the analysis of the upcoming Rubin data, but also morphology in other bands, from radio to UV. We are open to hearing about other potential morphology tracers that we can test and add to the code!

\hypertarget{abstract-talk-321}{}\label{abstract-page-talk-321}
\subsubsection*{Spatially-resolved star formation in the Virgo cluster}
\addcontentsline{toc}{subsubsection}{Spatially-resolved star formation in the Virgo cluster}
\textbf{Speaker:} Cameron Morgan\\
\textbf{Fields:} Galaxy $\cdot$ Galaxy Clusters

The Virgo cluster presents a unique opportunity to disentangle the roles of environmental quenching mechanisms such as ram-pressure stripping and starvation given its proximity and ongoing formation. Leveraging deep, spatially resolved optical and Halpha data from Virgo cluster surveys, we are probing the effects of the environment on cluster galaxies on integrated and spatially resolved scales. We have measured stellar mass functions across the cluster, splitting the population according to star formation rate. The results of this work show that there is little change in the shape of the mass functions across the cluster, and quenched fractions are particularly high for dwarf galaxies, even beyond the cluster virial radius. As a first morphological indicator of quenching, we have quantified the edges of star-forming disks and shown that galaxies in the cluster exhibit truncations with respect to normal star-forming galaxies. The ubiquity of this feature across the cluster, including in areas where ram-pressure stripping cannot be dominant, suggests that more passive mechanisms such as starvation consume gas and kickstart quenching well before the galaxy approaches first pericenter. To extend this study beyond these initial tests, we are working to constrain spatially resolved star formation histories in Virgo galaxies that have different environmental quenching signatures, to understand how recent star formation is affected by various quenching mechanisms. I will share exciting new results from this final study, which highlight the power of spatially-resolved data to aid in our understanding of galaxy evolution.

\hypertarget{abstract-talk-19}{}\label{abstract-page-talk-19}
\subsubsection*{Horizon-scale dynamics in the accretion flow of the supermassive black hole in Sagittarius A*}
\addcontentsline{toc}{subsubsection}{Horizon-scale dynamics in the accretion flow of the supermassive black hole in Sagittarius A*}
\textbf{Speaker:} Rohan Dahale\\
\textbf{Field:} Supermassive black hole

The Event Horizon Telescope (EHT) has imaged the time-averaged structure of black hole at the Galactic Center, Sagittarius A* (Sgr A*). Resolving its temporal evolution enables new science: directly observing how matter accretes and how plasma evolves on minute timescales near the event horizon, testing accretion physics beyond static images. However, Sgr A*'s rapid variability and sparse interferometric coverage make dynamic imaging challenging. We now aim to spatially resolve the time evolution of Sgr A* by reconstructing a minute-by-minute horizon scale video from the April 11, 2017 EHT observations, the best dataset of the 2017 campaign, with $\sim$3 hours of continuous coverage and ten baselines. This day showed pronounced multi-wavelength variability: an X-ray flare detected by Chandra X-ray Observatory and NuSTAR immediately before the EHT track; Q-U loops in polarimetric light curves from ALMA; and quasi-periodic oscillations in EHT polarization visibilities, all suggesting coherent dynamics in the accretion flow. To extract robust videos, we developed (1) independent dynamic imaging methods using neural fields, Bayesian inference, and optimal-transport regularization, and (2) a validation framework testing whether these methods distinguish genuine motion from artifacts, using synthetic data from geometric models to GRMHD simulations. The validated reconstructions will reveal horizon-scale accretion dynamics and enable direct measurements of rotational plasma velocities and evolving brightness structure. Combined with multi-wavelength data following the flare, they constrain magnetic-field geometry, plasma properties, and flare mechanisms, opening a new frontier for understanding variability in low-luminosity black hole accretion and testing strong-gravity predictions.

\hypertarget{abstract-talk-145}{}\label{abstract-page-talk-145}
\subsubsection*{Cold and warm molecular gas in MACS1931-26 brightest cluster galaxy}
\addcontentsline{toc}{subsubsection}{Cold and warm molecular gas in MACS1931-26 brightest cluster galaxy}
\textbf{Speaker:} Laya Ghodsi\\
\textbf{Fields:} Galaxy $\cdot$ Galaxy Clusters

Galactic star formation, and consequently the evolutionary path of galaxies, are shaped by their gas content within the interstellar medium (ISM), which directly fuels the formation of stars. The ISM is influenced by the interactions of galaxies with their close environment, the circumgalactic medium (CGM), a diffuse gas reservoir surrounding them. The CGM acts as a bridge between ISM of galaxies and their cosmological environment, enabling gas flows that fuel star formation and redistribute metals throughout the medium. I will discuss the multi-temperature molecular ISM and CGM of the brightest cluster galaxy (BCG) in the cool-core cluster MACS1931-26 at z=0.35. This galaxy hosts one of the largest known H$_{2}$ reservoirs, exhibits elevated star formation, and contains a radio-loud AGN. We trace cold H$_{2}$ (10-100 K) in MACS1931-26 using multiple CO and CI emission lines, revealing that the gas in the ISM is highly excited, comparable to local ULIRGs. In contrast, the CGM gas is much less excited, indicating differing ionization sources in these environments. Our JWST observations reveal warm H$_{2}$ (100-1000 K) spatially coincident with CO with similar kinematics. Our analysis of the H2 rotational lines shows a higher excitation temperature for the CGM compared to the ISM. The CGM also has a higher warm-to-cold H2 mass fraction than the ISM. In this system, both shocks and AGN-driven X-ray emission contribute to heating the H$_{2}$-emitting gas. Ongoing shock modeling aims to reveal potential gas flows from the CGM to the ISM of this galaxy, which may fuel star formation and correlate with the strong activity of AGN, shedding light on the complex feedback processes governing galaxy evolution.

\hypertarget{abstract-talk-51}{}\label{abstract-page-talk-51}
\subsubsection*{Two Wavelengths, Two Components: A New View of Galaxy Growth}
\addcontentsline{toc}{subsubsection}{Two Wavelengths, Two Components: A New View of Galaxy Growth}
\textbf{Speaker:} Angelo George\\
\textbf{Field:} Galaxy

Galaxy structural evolution is encoded in the sizes of their bulges and disks. Using deep imaging from the CFHT CLAUDS and Subaru HSC-SSP surveys, we analyze a mass-complete sample of $\sim$44,000 star-forming galaxies and $\sim$34,000 quiescent galaxies and perform bulge+disk decompositions in two rest-frame wavelengths (UV at 3000 and optical at 5000 ). This enables the first large statistical dual-wavelength analysis of the sizemass relation and its evolution for bulges and disks separately.

\hypertarget{abstract-talk-327}{}\label{abstract-page-talk-327}
\subsubsection*{Enhanced AGN fraction in the protocluster SPT2349-56}
\addcontentsline{toc}{subsubsection}{Enhanced AGN fraction in the protocluster SPT2349-56}
\textbf{Speaker:} Vismaya Pillai\\
\textbf{Fields:} Cosmology and Early Universe $\cdot$ Galaxy $\cdot$ Galaxy Clusters

Active Galactic Nuclei are important drivers of galaxy evolution through both positive and negative feedback mechanisms. The dense environment of high-redshift clusters and protoclusters is conducive to high rates of both star formation and SMBH accretion. The AGN fraction of a protocluster is thus an important predictor of cluster evolution.

We find an overabundance of obscured AGN in the extreme z=4.3 protocluster SPT2349-56 identified using the observed excess in rest-frame 3 $\mu$m flux. We perform CIGALE SED-fitting of new JWST MIRI and NIRCam photometry combined with archival coverage from optical to submm wavelengths to determine the fraction of AGN emission in protocluster members. We find that the observed MIRI 18/10 $\mu$m flux excess traces this AGN fraction with a monotonic relation and low scatter. Approximately half of the protocluster members exhibit a significant AGN component, a substantial enhancement over the 10\% AGN fraction seen in field galaxies. This enhancement is likely driven by the dense protocluster environment due to accretion of infalling gas from the cosmic web fuelling the AGN as well as boosted activity resulting from galaxy interactions.

With such a high fraction of AGN, the impact of AGN feedback is liable to outweigh stellar feedback effects and contribute substantially to quenching cluster galaxies and heating of the nascent intracluster medium. Thus, a heightened AGN fraction supports the expected track of SPT2349's evolution towards a massive, quenched Coma-like cluster in the late universe.

\hypertarget{abstract-talk-193}{}\label{abstract-page-talk-193}
\subsubsection*{Classifying Quasar Types Without a Spectrum}
\addcontentsline{toc}{subsubsection}{Classifying Quasar Types Without a Spectrum}
\textbf{Speaker:} Adrien Hélias\\
\textbf{Fields:} Astrostatistics, Astroinformatics, and AI / Machine Learning $\cdot$ Supermassive black hole

Classifying quasars is important because it helps us understand accretion regimes, black hole mass scaling, disk instabilities and feedback processes. Type 1 quasars offer a direct view of the accretion disk and the broad-line region, whereas Type 2 quasars are obscured by dust and only the narrow-line region is visible. This classification process is traditionally done with spectroscopy because of its robustness. However, it does not scale well with millions of quasars, as it requires substantial time and resources to acquire a good spectrum. In this work, we show that we can use irregularly-sampled photometry (light curves) to classify quasar types without a spectrum, using Slepian Wavelet Variance (SWV). We differentiate Type 1 and Type 2 quasars from the MILLIQUAS catalogue, solely based on their Zwicky Transient Facility optical light curves. SWV provides comparable or better performance with fewer assumptions than applying structure functions or the Damped Random Walk model on these time series. In the era of large, fast-cadence surveys like the Legacy Survey of Space and Time, SWV is computationally fast enough to be used on a large number of light curves and offers promise for efficient classification of quasar type without spectroscopy.


\section*{Afternoon}
\addcontentsline{toc}{section}{Thursday -- Afternoon}

\subsection*{A-1502}
\textbf{Theme:} Galaxy Clusters

\hypertarget{abstract-talk-273}{}\label{abstract-page-talk-273}
\subsubsection*{Revealing Gas Inflow Toward the Supermassive Black Hole in NGC 4696 with the James Webb Space Telescope}
\addcontentsline{toc}{subsubsection}{Revealing Gas Inflow Toward the Supermassive Black Hole in NGC 4696 with the James Webb Space Telescope}
\textbf{Speaker:} Mathieu Marquis\\
\textbf{Fields:} Galaxy $\cdot$ Galaxy Clusters $\cdot$ Supermassive black hole

Understanding how supermassive black holes influence the galaxies they inhabit remains a central challenge in astrophysics. The most massive galaxies in the UniverseBrightest Cluster Galaxies at the centers of galaxy clustersoffer an exceptional laboratory to study this connection. Immersed in vast reservoirs of hot gas, their central black holes both accrete surrounding material and release enormous amounts of energy, profoundly shaping the galaxy and its environment. The nearby Centaurus cluster, whose central galaxy is NGC 4696, provides a unique opportunity to observe this interplay in action thanks to its proximity and its spectacular network of cold gas filaments.

In this talk, I will present new JWST/NIRSpec IFU observations that, for the first time, allow us to follow the entire feedback cycle linking black holes and galaxiesfrom gas cooling and condensation to accretion onto the black holeacross scales ranging from kpc down to pc. These observations reveal a rotating circumnuclear disk in NGC 4696 that is directly connected to the surrounding cold filament network. This structure provides the long-sought missing link that channels gas from kpc scales down to $\sim$100 pc and ultimately into the central black hole, effectively closing the feedback loop. Comparisons with stellar kinematics further show that gas and stars respond very differently to the central gravitational potential and to feedback processes. Finally, by combining JWST with MUSE, ALMA, and XRISM, we provide a multiwavelength view of the gas in the galaxy's core and its motions across scales.

These results provide new insight into how supermassive black holes shape their host galaxies and illustrate the transformative power of JWST for resolving the physical processes driving galaxy evolution.

\hypertarget{abstract-talk-29}{}\label{abstract-page-talk-29}
\subsubsection*{Forged by Feedback: How AGN Feedback Models Shape Brightest Group Galaxies in Cosmological Simulations}
\addcontentsline{toc}{subsubsection}{Forged by Feedback: How AGN Feedback Models Shape Brightest Group Galaxies in Cosmological Simulations}
\textbf{Speaker:} Ruxin Barré\\
\textbf{Fields:} Galaxy $\cdot$ Galaxy Clusters $\cdot$ Supermassive black hole

Most modern galaxy formation models can successfully reproduce galaxy population statistics, yet many still struggle with capturing the properties of massive galaxies. These discrepancies often arise from the modelling of active galactic nucleus (AGN) feedback and its regulation of gas cooling and star formation. It is therefore crucial to refine the realism of AGN feedback models in cosmological simulations. Brightest group galaxies (BGGs) at the centres of galaxy groups are shaped by AGN-driven cycles of heating and cooling, which provides an ideal laboratory for testing feedback models. We use BGGs as probes of feedback, investigating their properties across multiple cosmological simulations with distinct AGN feedback models: thermal feedback in ROMULUS, differing calibrations of two-regime kinetic feedback in SIMBA and SIMBA-C, and the novel, three-regime OBSIDIAN model, which maps accretion flow geometry to AGN-driven outflows. Our results reveal that the properties and quenching pathways of BGGs are highly sensitive to the AGN feedback model. We contrast these predictions against COSMOS BGGs and find that OBSIDIAN achieves the highest agreement with observations, reproducing BGG stellar properties and the evolution of star formation and quenching with mass. We find evidence that OBSIDIAN and COSMOS BGGs quench via a gradual decline in star formation. This contrasts SIMBA and SIMBA-C, which exhibit rapid quenching linked to the onset of jet feedback, and ROMULUS, whose thermal feedback cannot suppress star formation. These results provide critical constraints for refining AGN feedback models, with OBSIDIAN's success signalling the start of a new generation of simulations with increasingly realistic depictions of AGN and their impact on galaxy evolution.

\hypertarget{abstract-talk-144}{}\label{abstract-page-talk-144}
\subsubsection*{Mapping AGN Driven Outflows in the Local Universe with MaNGA}
\addcontentsline{toc}{subsubsection}{Mapping AGN Driven Outflows in the Local Universe with MaNGA}
\textbf{Speaker:} Sana Momin\\
\textbf{Fields:} Galaxy $\cdot$ Supermassive black hole

AGN driven outflows are considered an important mechanism in galaxy evolution, yet their spatial and dynamical impact remain uncertain. We present a spatially resolved study of ionized gas outflows in nearby AGN using integral field spectroscopy from the MaNGA survey. Asymmetries in emission line profiles serve as tracers of outflowing gas. Using the [O III] $\lambda$5007 line, we trace AGN driven outflows by identifying emission line asymmetries and isolating the kinematic outflow component across each galaxy. The spatially resolved MaNGA data enable us to locate where these asymmetric signatures arise and characterize key outflow properties, including their spatial extent, morphology, and kinematic strength. Using this approach, we generate two dimensional maps of outflow kinematics and projected radii, which allow us to directly compare outflow strength to AGN luminosity. Using the Comerford et al. (2020) MaNGA AGN catalog, we identify 33 AGN host galaxies exhibiting distinct outflow signatures within their AGN ionized regions. The outflows show a wide range of morphologies, from compact structures to extended features consistent with orientation-dependent effects. On average, the outflows reach a projected radii of 3.6 kpc, with mean outflow velocities of $\sim$202 km/s and velocity dispersions of 380 km/s. These results demonstrate that AGN driven outflows commonly extend to kpc scales and provide quantitative constraints on the strength and structure of AGN feedback.

\subsection*{A-3561}
\textbf{Theme:} Education and Public Outreach

\hypertarget{abstract-talk-180}{}\label{abstract-page-talk-180}
\subsubsection*{Tactile Galaxy Models for Accessible Astronomy}
\addcontentsline{toc}{subsubsection}{Tactile Galaxy Models for Accessible Astronomy}
\textbf{Speaker:} Grant Hurley\\
\textbf{Fields:} Education and Public Outreach $\cdot$ Stars and Stellar cluster formation

Astronomical data are typically communicated through visual representations such as images, plots, and maps. While effective for many audiences, these approaches can present barriers for learners with visual impairments and can make it challenging for students to fully grasp the three-dimensional structure of astronomical systems. Physical models provide an alternative way to explore spatial information and may help bridge this gap.

In this work I present a series of three-dimensional printed galaxy models generated from astronomical imaging and survey data. These models are designed to highlight key structural features such as spiral arms, central bulges, and extended disks while incorporating textures that allow features to be explored through touch. The models are being developed for use in classrooms, outreach activities, and museum-style displays to support tactile learning in astronomy.

By transforming observational data into tangible objects, these models provide an accessible way to engage with galaxy structure and morphology. This work explores the role that tactile models can play in astronomy education and demonstrates how emerging fabrication tools can be used to make astronomical data more inclusive and accessible to a broader range of learners.

\hypertarget{abstract-talk-55}{}\label{abstract-page-talk-55}
\subsubsection*{Beading Astroparticle Physics: Collaborating with the A.B. McDonald Canadian Astroparticle Physics Research Institute to translate astrophysical knowledge in beadwork}
\addcontentsline{toc}{subsubsection}{Beading Astroparticle Physics: Collaborating with the A.B. McDonald Canadian Astroparticle Physics Research Institute to translate astrophysical knowledge in beadwork}
\textbf{Speaker:} Danielle Lussier\\
\textbf{Fields:} Education and Public Outreach $\cdot$ Equity, Diversity and Inclusion $\cdot$ Indigenous Engagement $\cdot$ Particle Astrophysics

In collaboration with the Arthur B. McDonald Canadian Astroparticle Physics Research Institute Dr. Lussier (Red River Mtis legal scholar) and Dr. Wade (settler astrophysicist) are developing beadwork-based artwork as vehicles for astrophysical knowledge translation.

Beading works as a science methodology because it is both a material practice and a knowledge framework. The act of beading involves precision, pattern recognition, spatial reasoning, and iterative design, which are all skills that resonate with scientific processes. At the same time, beadwork is deeply embedded in Indigenous epistemologies, carrying teachings about interconnection, relationality, and cosmology. As a methodology, it emphasizes embodied learning (knowledge produced through hands, senses, and movement) and relational accountability (placing the learner within networks of people, materials, and land). This makes it particularly powerful for astroparticle physics, as it challenges linear, extractive approaches by offering a way of doing research that is creative, relational, and holistic, opening new avenues for understanding complexity, systems, and interdependence.

In this presentation we motivate the project, describe the objectives, and provide the CASCA community with a sneak peek into the knowledge translation works in progress.

\hypertarget{abstract-talk-11}{}\label{abstract-page-talk-11}
\subsubsection*{Introducing loom beadwork into an introductory astronomy class}
\addcontentsline{toc}{subsubsection}{Introducing loom beadwork into an introductory astronomy class}
\textbf{Speaker:} Gregg Wade\\
\textbf{Field:} Education and Public Outreach

We describe the introduction of loom beadwork into a small introductory astronomy course for Science students at Royal Military College of Canada. Inspired by the widespread beading traditions of Indigenous peoples across Canada, we use looms and glass seed beads to encourage learners to translate astronomical knowledge into aesthetically compelling artwork imbued with layers of knowledge.

Beading is often a form of mnemonic translation of knowledge. Loom beading encodes information on a rectilinear grid, and in this sense it is analogous to digital astronomical imagery, with beads analogous to individual pixels having both spatial characteristics (location within a two-dimensional matrix) and intensity characteristics (colour). This close similarity opens conversations about image scale, spatial resolution, and information content in astronomical imagery.

Beading also prompts conversations about Indigenous peoples of Canada, including pre-contact social organization and technologies, post-contact historical and contemporary realities, the TRC and its Calls to Action, and reconciliation. As a consequence beading can naturally be introduced into, or can itself open, discussions of different global cultural perspectives and achievements in astronomy.

In this presentation we will describe our motivation, process, and experience, and discuss lessons learned and plans moving forward.

\hypertarget{abstract-talk-211}{}\label{abstract-page-talk-211}
\subsubsection*{Astro-Bubble on Tour: Planetarium shows for rural communities}
\addcontentsline{toc}{subsubsection}{Astro-Bubble on Tour: Planetarium shows for rural communities}
\textbf{Speaker:} Roan Haggar\\
\textbf{Field:} Education and Public Outreach

Don't you hate it when you really want to visit a planetarium, but don't have the time to drive for hours to reach the closest one? The Astro-Bubble is a portable, inflatable planetarium, operated by the Waterloo Centre for Astrophysics at the University of Waterloo; it has allowed over 11,000 people to take part in nearly 500 free in-person planetarium shows in public spaces since late-2022, learning about topics ranging from stargazing and solar eclipses, to JWST and galaxy evolution.

Astro-Bubble planetarium shows take place in venues such as schools, community centres, and libraries in the Waterloo Region, as well as on the University of Waterloo campus. Accessibility to all is one of the main priorities for these shows, but like many science outreach initiatives run by academic institutions, access is naturally biased towards people living nearby to the University.

With this in mind, in October 2025 we took the Astro-Bubble on a week-long tour around Northern Ontario, visiting five rural schools and a community sports centre, to provide astronomy-based educational experiences to communities that might otherwise struggle to access them. We ran planetarium shows for 1,000 elementary students, high school students, and adult residents of these communities across the week, speaking to students about some of the biggest unanswered questions in astrophysics, as well as hearing their thoughts on space and astronomy. We will share some of our experiences, lessons learned, and plans for the future, based on this week-long science outreach activity.

\subsection*{A-4502}
\textbf{Theme:} Night-Sky Protection / Solar System

\hypertarget{abstract-talk-319}{}\label{abstract-page-talk-319}
\subsubsection*{The Luminosity Function Of Ultra-Faint Trans-Neptunian Objects Detected By JWST}
\addcontentsline{toc}{subsubsection}{The Luminosity Function Of Ultra-Faint Trans-Neptunian Objects Detected By JWST}
\textbf{Speaker:} Marielle Eduardo\\
\textbf{Field:} Solar System

The size distribution (SD) of Trans-Neptunian Objects (TNOs) has remained relatively unaltered since the formation of the Solar System. This makes it a key test for planet formation models, as it reflects the outcomes of accretion, collisions, and dynamical evolution of the Solar System. Comparing the observed SD with those predicted by models helps to constrain the physical processes and underlying initial conditions that shaped the current Solar System. However, ground-based surveys are limited to only about m(r)$\sim$27 mag, and the lack of observations of smaller TNOs (m(r)$\sim$28, D$\sim$20km) leaves models poorly constrained.

Using images obtained from JWST Cycle 1 GO Program \#1568, we detected and characterized TNOs as faint as m(r)$\sim$29.8 mag and as small as $\sim$7 km in diameter to explore never-before probed regions of the TNO SD. To identify moving objects efficiently, we employed shift-and-stack technique and a machine learning model trained to identify small moving bodies in JWST images. This is by far the deepest Solar System survey to date, with at least a visible magnitude deeper than the landmark survey by Bernstein et al. (2004) that used the HST.

Our results show that hot and cold TNO populations have similarly shallow SD slopes down to $\sim$7 km, consistent with planetesimal formation via streaming instability followed by pebble cloud collapse (SI+PCC). The similar slopes suggest both populations formed through comparable processes despite different birth regions in the protoplanetary disk, implying that planetesimal formation is largely insensitive to local disk conditions. The slightly shallower observed slope compared to model predictions hints at a characteristic mass scale in planetesimal formation and provides constraints for leading models.

\hypertarget{abstract-talk-244}{}\label{abstract-page-talk-244}
\subsubsection*{Research Announcements For The Solar System (RAFTS)}
\addcontentsline{toc}{subsubsection}{Research Announcements For The Solar System (RAFTS)}
\textbf{Speaker:} Laura Buchanan\\
\textbf{Field:} Solar System

With the Legacy Survey of Space and Time (LSST) Solar System Collaboration (SSSC), and the Canadian Astronomy Data Centre (CADC) we have created a research announcement system for the solar system. Named the Research Announcements for the Solar System (RAFTS), this is a publication system designed to quickly issue short, citable announcements relevant to solar system science. This system is intended to address the unique needs of the solar system research community in the era of large-scale surveys (such as LSST), providing a platform for timely dissemination of findings. Hosted by the CADC, RAFTS is freely accessible to all communities, and easily discoverable with permanent DOIs assigned to each announcement.

Each announcement will feature a machine-readable section to facilitate rapid follow-up, and integration with the LSST community forum will encourage further collaborations, discussion and engagement. The system also features a moderation process to ensure high-quality and relevant publications. The announcements are expected to be brief, relevant, and may be urgent, though urgency is not a requirement for publication. The system is designed to be scalable, with the capacity to handle an increased volume of discoveries expected from the LSST. We will present a brief demonstration of the software interface, highlighting its user-friendly design and functionality for the community.

\hypertarget{abstract-talk-305}{}\label{abstract-page-talk-305}
\subsubsection*{Things that flash and beep in the night: future technologies and the changing sky}
\addcontentsline{toc}{subsubsection}{Things that flash and beep in the night: future technologies and the changing sky}
\textbf{Speaker:} Aaron Boley\\
\textbf{Field:} Night

Space mirrors, orbital data centres, satellite refuelling depots, meteor showers on demand, space-based missile defence, and obtrusive space-based advertising are all examples of plans that are currently being pursued that could transform low Earth orbit and the night sky. While a sky full of 100,000 satellites was once seen by some as hyperbolic, powerful actors throughout the world are rushing to make such a situation a reality. It is not just numbers that are concerning, but also the use cases. One example is provisioning daylight as a service, where satellites are designed to purposefully redirect sunlight onto Earth for various purposes. A prototype, which could launch this year, is expected to achieve illumination that is equivalent to several times the brightness of the full Moon, possibly exceeding a similar test conducted in the 1990s. As another example, some systems are being put forward as a way to save the environment by launching computational facilities into space. One specific case, filed at a national regulator, would blanket local skies with visible satellites, possibly exceeding the total number of stars visible by the unaided eye. Here, we provide an overview of some of these plans, potential impacts, and the international response so far. These impacts go beyond dark and quiet skies, and include concerns over rocket exhaust and satellite reentry emissions, space debris, and orbital collision and reentry risks. We discuss ongoing international mechanisms and options for addressing these concerns, as well as actions that the Canadian astronomical community can take to help preserve the sky.

\hypertarget{abstract-talk-177}{}\label{abstract-page-talk-177}
\subsubsection*{Catastrophic disruption of asteroid 2023 CX1 and implications for planetary defence}
\addcontentsline{toc}{subsubsection}{Catastrophic disruption of asteroid 2023 CX1 and implications for planetary defence}
\textbf{Speaker:} Auriane Egal\\
\textbf{Field:} Solar System

On February 13, 2023, asteroid 2023 CX1 entered Earths atmosphere over Normandy, producing a spectacular fireball and the fall of the Saint-Pierre-le-Viger meteorites. Discovered only seven hours before impact, it became the first asteroid ever studied in space, in the atmosphere, and in the laboratory. Thanks to an unprecedented international collaboration, fragments were recovered within hours; a success made possible by the most precise trajectory prediction ever obtained for an impactor and the first large-scale, targeted observation of a fireball, supported by exceptional public mobilization. Working together, astronomers, citizen scientists, and local communities retraced for the first time the complete history of an asteroid, from its origin in the main belt to its laboratory analysis. The recovered meteorites, ordinary L chondrites, are among the most common on Earth yet they revealed an unexpected behavior. Instead of breaking up gradually in the atmosphere, 2023 CX1 exploded violently at 28 km altitude, releasing nearly all its energy at once. This abrupt fragmentation produced a concentrated shockwave, posing greater risks to populated areas than the more progressive disintegrations usually observed. Such findings have direct implications for planetary defense, highlighting the importance of rapid detection, accurate prediction, and coordinated response.
